\chapter{Egg implementation of the saturation example}\label{appendix:smce-rules}

This appendix documents the computational counterpart of~\autoref{example:sat-comonoid} using the $\egg$ equality saturation library~\cite{EggPaper}.
We describe how the symmetric monoidal theory of commutative comonoids is encoded in the $\egg$ framework and explain the design decisions that shape the implementation.

In $\egg$, the term language is a Rust \mintinline{rust}{enum} defined via the \mintinline{rust}{define_language!} macro, which associates string tokens with constructors and fixes their arities.
\cref{lst:language} shows the language for commutative comonoids.

\begin{listing}[h]
\caption{Term language for commutative comonoids.}
\label{lst:language}
\begin{minted}{rust}
define_language! {
    enum CommutativeComonoids {
        "mu"       = Mu,              // mu : 1 -> 2
        "i"        = I,               // i  : 1 -> 0  (counit eta)
        ";"        = Compose([Id; 2]),
        "parallel" = Parallel([Id; 2]),
        "id_0" = Identity0,  "id_1" = Identity1,
        "id_2" = Identity2,  "id_3" = Identity3,  "id_4" = Identity4,
        "sym_0_1" = Sym0_1,  "sym_1_0" = Sym1_0,  "sym_1_1" = Sym1_1,
        "sym_1_2" = Sym1_2,  "sym_2_1" = Sym2_1,  "sym_2_2" = Sym2_2,
    }
}
\end{minted}
\end{listing}

This particular encoding is influenced by the following considerations.

The SMC equations reference identities $\id_n$ and symmetries $\sym_{n,m}$ for all $n,m\in\mathbb{N}$.
Because $\egg$ requires a finite language, these must include the arities actually needed by the example: identities at arities $0$--$4$ and symmetries for all $(n,m)$ with $n+m \leq 3$ and additionally $\sym_{2,2}$.
This necessitates the \emph{definitional rules} (e.g.\ $\id_2 = \id_1 \otimes \id_1$, $\sym_{2,1} = (\sym_{1,1} \otimes \id_1);(\id_1 \otimes \sym_{1,1})$) that link these tokens.
Each definitional rule is applied at most once during saturation and is omitted from~\cref{tab:appendix-rules}.

String diagrams allow $n$-ary parallel composition, but $\egg$ operators have fixed binary arity.
Encoding $\otimes$ as a binary \mintinline{rust}{"parallel"} operator requires explicit \emph{tensor associativity} rules $(f \otimes g) \otimes h = f \otimes (g \otimes h)$ and \emph{tensor unit} laws $\id_0 \otimes f = f$, $f \otimes \id_0 = f$ to expose all equivalent forms of a tensor expression.
As~\cref{tab:appendix-rules} shows, these two rule groups account for more than half of all applications---a direct cost of the encoding.

Each morphism in an SMC has a type $(dom, cod)\in\mathbb{N}^2$.
The $\egg$ framework supports \emph{e-class analyses}: user-defined data associated with each equivalence class, maintained invariantly as classes are merged.
A \mintinline{rust}{TypeAnalysis} computes the type of every expression bottom-up, as shown in~\cref{lst:types}.

\begin{listing}[h]
\caption{Type computation for composition and tensor product.}
\label{lst:types}
\begin{minted}{rust}
Compose([f, g]) => Ty { dom: get(f).dom, cod: get(g).cod },
Parallel([f, g]) => Ty {
    dom: get(f).dom + get(g).dom,
    cod: get(f).cod + get(g).cod,
},
\end{minted}
\end{listing}

Types are used to \emph{guard} rules that are only valid at specific arities.
For example, the left-identity rule $\id_n ; f = f$ must be restricted to morphisms $f$ with domain $n$:

\begin{listing}[h]
\caption{Type-guarded identity rules for $n=1$ (the pattern repeats for $n=2,3,4$).}
\label{lst:identity}
\begin{minted}{rust}
rw!("id_1-left-compose-left";  "(; id_1 ?x)" => "(?x)"
    if check_dom_type(var("?x"), 1)),
rw!("id_1-left-compose-right"; "(?x)" => "(; id_1 ?x)"
    if check_dom_type(var("?x"), 1)),
\end{minted}
\end{listing}

Without the guard, the rule would fire for morphisms of any type, producing ill-typed terms and incorrect equalities.
The guard is evaluated against the \mintinline{rust}{TypeAnalysis} data stored in each e-class.

\cref{lst:rules} shows the remaining representative rules.
Each equation of the theory is split into two directed rules, one for each orientation, so that equality saturation can add both sides of the equivalence to the same e-class.

\begin{listing}[h]
\caption{Representative rewrite rules (simplified for presentation).}
\label{lst:rules}
\begin{minted}{rust}
// Composition associativity
rw!("assoc-compose-left";  "(; ?x (; ?y ?z))" => "(; (; ?x ?y) ?z)"),
rw!("assoc-compose-right"; "(; (; ?x ?y) ?z)" => "(; ?x (; ?y ?z))"),

// Tensor associativity
rw!("parallel-assoc-left";
    "(parallel ?x (parallel ?y ?z))" => "(parallel (parallel ?x ?y) ?z)"),
rw!("parallel-assoc-right";
    "(parallel (parallel ?x ?y) ?z)" => "(parallel ?x (parallel ?y ?z))"),

// Symmetry naturality for cod(f)=cod(g)=1
// (general form: sym_{1,1};(f*g) = (g*f);sym_{cod(g),cod(f)})
rw!("sym-nat-1-1";
    "(; sym_1_1 (parallel ?f ?g))" => "(; (parallel ?g ?f) sym_1_1)"
    if check_dom_cod(1,1, 1,1)),  // dom and cod of f and g

// Comonoid laws
rw!("comult-commutative-right"; "(; mu sym_1_1)" => "(mu)"),
rw!("comult-commutative-left";  "(mu)" => "(; mu sym_1_1)"),
rw!("counit-left"; "(; mu (parallel i id_1))" => "(id_1)"),
\end{minted}
\end{listing}

The symmetry naturality rule $\sym_{1,1};(f \otimes g) = (g \otimes f);\sym_{cod(g),cod(f)}$ cannot be expressed as a single rule because the output symmetry token depends on the codomains of $f$ and $g$, which are run-time data rather than syntactic structure.
The rule is therefore split into one instance for each possible combination of codomains, each guarded by type conditions on $f$ and $g$.
Six such instances are required for the example, covering codomains $(cod(f), cod(g)) \in \{(1,1),(1,0),(0,1),(1,2),(2,1),(2,2)\}$.

The initial e-graph contains the term
\[
  ((\mu;\sym_{1,1});(\eta \otimes \id_1)) \otimes (\mu;(\id_1 \otimes \id_1)),
\]
represented in the $\egg$ S-expression syntax as
\begin{minted}{text}
(parallel (; (; mu sym_1_1) (parallel i id_1))
          (; mu (parallel id_1 id_1)))
\end{minted}
After saturation, the code checks that this e-class equals \mintinline{text}{(parallel id_1 mu)}, i.e.\ $\id_1 \otimes \mu$, returning \mintinline{rust}{true}.
The equivalence follows by: commutativity ($\mu;\sym_{1,1} = \mu$) followed by the counit law ($\mu;(\eta \otimes \id_1) = \id_1$) for the left component, and the right identity law ($\mu;(\id_1 \otimes \id_1) = \mu;\id_2 = \mu$) for the right component.

\Cref{tab:appendix-rules} reports the number of times each non-definitional rule was applied during saturation.
As we already noted in~\cref{chapter:introduction}, the domain specific rules of commutativity and counit are applied exactly once, whereas structural rules are applied hundreds of times.

\begin{table}[h]
\centering
\begin{tabular}{@{}llp{7.5cm}rr@{}}
\toprule
Group & Rule & Equation & $(\to)$ & $(\leftarrow)$ \\
\midrule
\multirow{8}{*}{SMC}
  & Left identity       & $\id_n ; f = f$                                                                     & ---     & 899 \\
  & Right identity      & $f ; \id_n = f$                                                                     & 4       & 308 \\
  & Composition assoc.  & $(f ; g) ; h = f ; (g ; h)$                                                         & 1{,}325 & 580 \\
  & Bifunctoriality     & $(f_1 \otimes f_2) ; (g_1 \otimes g_2) = (f_1 ; g_1) \otimes (f_2 ; g_2)$           & 157     & 92  \\
  & Tensor assoc.       & $(f \otimes g) \otimes h = f \otimes (g \otimes h)$                                  & 1{,}695 & 2{,}174 \\
  & Tensor unit         & $\id_0 \otimes f = f,\quad f \otimes \id_0 = f$                                     & ---     & 2{,}005 \\
  & Sym.\ naturality    & $\sym_{1,1} ; (f \otimes g) = (g \otimes f) ; \sym$                                 & 10      & ---   \\
  & Sym.\ involution    & $\sym ; \sym = \id$                                                                  & ---     & 1     \\
\midrule
\multirow{2}{*}{$(\Sigma,\mathcal{E})$}
  & Commutativity       & $\mu ; \sym = \mu$                                                                   & 1       & --- \\
  & Counit              & $\mu ; (\eta \otimes \id) = \id$                                                     & 1       & --- \\
\midrule
  & \multicolumn{2}{l}{Total (shown rules only)} & 3{,}193 & 6{,}059 \\
\bottomrule
\end{tabular}
\caption{Rewrite rules applied during saturation in~\cref{example:sat-comonoid}. The subscript $n$ in the identity rules ranges over $\{1,2,3,4\}$. The $(\to)$ and $(\leftarrow)$ columns give application counts for the left-to-right and right-to-left orientations of each equation; --- indicates zero applications.}
\label{tab:appendix-rules}
\end{table}

This result is unsurprising, the encoding we chose did not seem to be optimal in the first place.
A possible alternative is to use the internal language of symmetric\footnote{Technically, these rules are for the linear monoidal categories, the symmetric version has additional conditions on the contexts.} monoidal categories~\cite{shulman2016categorical} given as follows:

\[
\inferrule*[right=$id$]{\vdash A\;\text{type}}{x : A \vdash x: A}
\]
\[
\inferrule*[right=$f_{\text{intro}}$]{f \in \mathcal{G}(A_1, \ldots, A_n; B)\;\Gamma_1 \vdash M_1 : A_1, \ldots, \Gamma_n \vdash M_n : A_n}{\Gamma_1, \ldots, \Gamma_n \vdash f(M_1, \ldots, M_n) : B}
\]
\[
\inferrule*[right=$\otimes_{\text{intro}}$]{\Gamma \vdash N : A \qquad \Delta \vdash M : B}{\Gamma, \Delta \vdash \{N \otimes M\} : A \otimes B}
\]
\[
\inferrule*[right=$\otimes_{\text{elim}}$]{\Psi \vdash M : A \otimes B \quad \Gamma, x : A, y : B, \Delta \vdash N : C}{\Gamma, \Psi, \Delta \vdash \text{match}_{A \otimes B}(M, xy.N) : C}
\]
\[
\inferrule*[right=$\I_{\text{intro}}$]{\;}{() \vdash \star : \I} \qquad \inferrule*[right=$\I_{\text{elim}}$]{\Psi \vdash M : \I \qquad \Gamma,\Delta \vdash N : A}{\Gamma, \Psi, \Delta \vdash \text{match}_{\I}(M,N) : A}
\]

such that $\mathcal{G}$ contains the generators of the theory, i.e., $\mu : 1 \to 2$ and $\eta : 1 \to 0$ for the comonoid theory.

In this language, our initial term can be expressed as follows:
\[
  x : A, y : A \vdash
  \{ \mathrm{match}_{A \otimes A}(\mu(x), ab.\ \mathrm{match}_{\I}(\eta(b), a)) \otimes \mu(y) \}
  : A \otimes (A \otimes A)
\]
This representation is slightly more economical, in particular, it does not require the explicit encoding of the $\id$ morphsisms, or equations for associativity of composition and its unitality.
However, it requires a much more sophisticated encoding of the term language which does not seem to be supported by \egg.
Specifically, the variable contexts and variable binding are necessary.


% \paragraph{SMC equations.}
% \emph{Left} and \emph{right identity} express that $\id_n$ is a unit for sequential composition.
% The strong asymmetry (899 introductions vs.\ 0 eliminations for left identity; 308 vs.\ 4 for right identity) is a hallmark of equality saturation: the algorithm adds both sides of every equation, so it frequently introduces $\id_n ; f$ as equivalent to $f$ to enable subsequent pattern matches, but rarely needs to simplify $\id_n ; f$ away.
% \emph{Composition associativity} is applied 1,905 times in total; more applications go right-to-left (flattening left-nested compositions) than left-to-right, reflecting which normal form is more useful for downstream matches.
% \emph{Bifunctoriality} converts between interleaved and separated compositions.
% \emph{Tensor associativity} (3,869 total) and the \emph{tensor unit} (2,005 total) together account for $56\%$ of all applications, directly reflecting the overhead of the binary \mintinline{rust}{"parallel"} encoding described above.
% \emph{Symmetry naturality} fires in only one direction and only 10 times, at the specific arity combinations needed to expose the commutativity of $\mu$.
% \emph{Symmetry involution} ($\sym;\sym = \id$, applied once in the introducing direction) provides the converse that allows the saturation to close the loop on symmetry-related equivalences.

% \paragraph{Comonoid equations.}
% \emph{Commutativity} ($\mu;\sym = \mu$) and the \emph{counit law} ($\mu;(\eta \otimes \id) = \id$) are each applied exactly once, confirming that the verification terminates as soon as the two relevant equivalences are found.
% Coassociativity of $\mu$ ($\mu;(\id \otimes \mu) = \mu;(\mu \otimes \id)$) is present as a rule but not applied: the particular input term does not require it.

\chapter{Additional material for Chapter~\ref{chapter:enriched}}

\section{Change of base and opposite categories}
\label{sec:appendix-change-of-base-op}

\begin{proposition}
Let $\mathcal{C}$ be an ordinary category and let $\widetilde{F}_{\mathcal{V}}$ be a change-of-base $2$-functor induced by a symmetric strong monoidal functor $F_{\mathcal{V}} : \catname{Set} \to \mathcal{V}$.
Then, we have that $\widetilde{F}_{\mathcal{V}}(\mathcal{C})^{\mathrm{op}} = \widetilde{F}_{\mathcal{V}}(\mathcal{C}^{\mathrm{op}})$.
\end{proposition}
\begin{proof}
The categories agree on objects.
The hom-objects of $\widetilde{F}_{\mathcal{V}}(\mathcal{C}^{\mathrm{op}})$ are given by 
\[
(\widetilde{F}_{\mathcal{V}}(\mathcal{C}^{\mathrm{op}}))(A,B) = F_{\mathcal{V}}(\mathcal{C}^{\mathrm{op}}(A,B)) = F_{\mathcal{V}}(\mathcal{C}(B,A))~.
\]
The hom-objects of $\widetilde{F}_{\mathcal{V}}(\mathcal{C})^{\mathrm{op}}$ are given by 
\[
(\widetilde{F}_{\mathcal{V}}(\mathcal{C})^{\mathrm{op}})(A,B) = (\widetilde{F}_{\mathcal{V}}(\mathcal{C}))(B,A) = F_{\mathcal{V}}(\mathcal{C}(B,A))~.
\]
Therefore, they also agree on hom-objects.
We need to check if the composition and identities agree.
For any triple of objects $A,B,C$, we have that the following diagram commutes
\[\begin{tikzcd}[column sep=small]
	&& {(\widetilde{F}_{\mathcal{V}}(\mathcal{C})^{\mathrm{op}})(B,C) \otimes (\widetilde{F}_{\mathcal{V}}(\mathcal{C})^{\mathrm{op}})(A,B)} \\
	{(\widetilde{F}_{\mathcal{V}}(\mathcal{C})^{\mathrm{op}})(A,B) \otimes (\widetilde{F}_{\mathcal{V}}(\mathcal{C})^{\mathrm{op}})(B,C)} && {(\widetilde{F}_{\mathcal{V}}(\mathcal{C}^{\mathrm{op}}))(B,C) \otimes (\widetilde{F}_{\mathcal{V}}(\mathcal{C}^{\mathrm{op}}))(A,B)} \\
	{\widetilde{F}_{\mathcal{V}}(\mathcal{C})(B,A) \otimes \widetilde{F}_{\mathcal{V}}(\mathcal{C})(C,B)} && {F_{\mathcal{V}}(\mathcal{C}^{\mathrm{op}}(B,C)) \otimes F_{\mathcal{V}}(\mathcal{C}^{\mathrm{op}}(A,B))} \\
	{F_{\mathcal{V}}(\mathcal{C}(B,A)) \otimes F_{\mathcal{V}}(\mathcal{C}(C,B))} && {F_{\mathcal{V}}(\mathcal{C}^{\mathrm{op}}(B,C) \times \mathcal{C}^{\mathrm{op}}(A,B))} \\
	{F_{\mathcal{V}}(\mathcal{C}(B,A) \times \mathcal{C}(C,B))} && {F_{\mathcal{V}}(\mathcal{C}^{\mathrm{op}}(A,C))} \\
	{\widetilde{F}_{\mathcal{V}}(\mathcal{C}(C,A))} && {F_{\mathcal{V}}(\mathcal{C}(C,A))}
	\arrow["{\id \otimes \id}"', from=1-3, to=2-1]
	\arrow["{=}"{description}, draw=none, from=1-3, to=2-3]
	\arrow["{\sym_{\mathcal{V}}}"', from=2-1, to=3-1]
	\arrow["{=}"{description}, draw=none, from=2-3, to=3-3]
	\arrow["{=}"{description}, draw=none, from=3-1, to=4-1]
	\arrow["{\Phi_{\mathcal{C}^{\mathrm{op}}(B,C),\mathcal{C}^{\mathrm{op}}(A,B)}}", from=3-3, to=4-3]
	\arrow["{\Phi_{\mathcal{C}(B,A),\mathcal{C}(C,B)}}"', from=4-1, to=5-1]
	\arrow["{F_{\mathcal{V}}(\circ_{\mathcal{C}^{\mathrm{op}}})}", from=4-3, to=5-3]
	\arrow["{F_{\mathcal{V}}(\circ_{\mathcal{C}})}"', from=5-1, to=6-1]
	\arrow["{=}"{description}, draw=none, from=5-3, to=6-3]
	\arrow["{=}"{description}, draw=none, from=6-1, to=6-3]
\end{tikzcd}~.\]
First note that $\circ_{\mathcal{C}^{\mathrm{op}}} = \circ_{\mathcal{C}} \circ \sym_{\catname{Set}}$, thus the equality between the left and right paths amounts to 
\[
F_{\mathcal{V}}(\circ_{\mathcal{C}}) \circ F_{\mathcal{V}}(\sym_{\catname{Set}}) \circ \Phi_{\mathcal{C}^{\mathrm{op}}(B,C),\mathcal{C}^{\mathrm{op}}(A,B)} = F_{\mathcal{V}}(\circ_{\mathcal{C}}) \circ \Phi_{\mathcal{C}(B,A),\mathcal{C}(C,B)} \circ \sym_{\mathcal{V}}
\] 
and we have that 
\[
F_{\mathcal{V}}(\sym_{\catname{Set}}) \circ \Phi_{\mathcal{C}^{\mathrm{op}}(B,C),\mathcal{C}^{\mathrm{op}}(A,B)} = \Phi_{\mathcal{C}(B,A),\mathcal{C}(C,B)} \circ \sym_{\mathcal{V}}~,
\] 
because the functor $F_{\mathcal{V}}$ is symmetric monoidal, as established in~\cref{cor:cocomplete_monoidal_free}.
Identity maps agree by the definition of $\widetilde{F}_{\mathcal{V}}(\mathcal{C})^{\mathrm{op}}$.
\end{proof}

\chapter{Additional material for Chapter~\ref{chapter:string_diagrams}}

\section{Pointwise pushouts of hypergraphs}
\label{sec:appendix-pointwise-pushouts}

Let $\mathcal{G}_0$, $\mathcal{G}_1$, and $\mathcal{G}_2$ be hypergraphs, viewed as functors $\mathcal{G}_i : \catname{I} \to \mathbb{F}$.
Let $f : \mathcal{G}_0 \to \mathcal{G}_1$ and $g : \mathcal{G}_0 \to \mathcal{G}_2$ be hypergraph homomorphisms, i.e., natural transformations.
Their pushout $\mathcal{G}_1 \coprod_{\mathcal{G}_0} \mathcal{G}_2$ is computed pointwise in $\mathbb{F}^{\catname{I}}$: for each object $x$ of $\catname{I}$ it is the pushout
\[
(\mathcal{G}_1 \coprod_{\mathcal{G}_0} \mathcal{G}_2)(x) = \mathcal{G}_1(x) \coprod_{\mathcal{G}_0(x)} \mathcal{G}_2(x)
\]
in $\mathbb{F}$ of the span $\mathcal{G}_1(x) \xleftarrow{f_x} \mathcal{G}_0(x) \xrightarrow{g_x} \mathcal{G}_2(x)$, which is given by $\mathcal{G}_1(x) \coprod \mathcal{G}_2(x)$ quotiented by the least equivalence relation $\sim_x$ such that $\iota^x_1(z_1) \sim_x \iota^x_2(z_2)$ if there exists $w \in \mathcal{G}_0(x)$ with $z_1 = f_x(w)$ and $z_2 = g_x(w)$, where $\iota^x_1, \iota^x_2$ are the injections into $\mathcal{G}_1(x) \coprod \mathcal{G}_2(x)$; we denote this quotient $(\mathcal{G}_1(x) \coprod \mathcal{G}_2(x)) / {\sim_x}$.
Then, given an arrow $\sigma : (k,l) \to \star$, we get an arrow $\alpha \defeq (\mathcal{G}_1 \coprod \mathcal{G}_2)(\sigma)$, $\alpha : (\mathcal{G}_1 \coprod \mathcal{G}_2)((k,l)) \to (\mathcal{G}_1 \coprod \mathcal{G}_2)(\star)$ and then the arrow $\hat{\alpha} : (\mathcal{G}_1 \coprod \mathcal{G}_2)((k,l)) \to (\mathcal{G}_1 \coprod \mathcal{G}_2)(\star) / \sim_{\star}$ by post-composing with the quotient projection.
We then let $\tilde{\alpha}([w]) = \hat{\alpha}(w)$ for $[w] \in (\mathcal{G}_1 \coprod \mathcal{G}_2)((k,l)) / {\sim_{(k,l)}}$.
That is, the action of $\mathcal{G}_1 \coprod_{\mathcal{G}_0} \mathcal{G}_2$ on arrows is given by $\tilde{\alpha}$.
Suppose $w_1 \sim_{(k,l)} w_2$, then $\hat{\alpha}(w_1) = \hat{\alpha}(w_2)$, that is, $\tilde{\alpha}$ is well-defined.
To illustrate that, assume there exists $z \in \mathcal{G}_0((k,l))$ such that $w_1 = \iota^{(k,l)}_1 \circ f_{(k,l)}(z)$ and $w_2 = \iota^{(k,l)}_2 \circ g_{(k,l)}(z)$, then we need to check if
\begin{gather*}
\hat{\alpha}(\iota^{(k,l)}_1 \circ f_{(k,l)}(z)) = \hat{\alpha}(\iota^{(k,l)}_2 \circ g_{(k,l)}(z))\\
\iff\\
\alpha(\iota^{(k,l)}_1 \circ f_{(k,l)}(z)) \sim_{\star} \alpha(\iota^{(k,l)}_2 \circ g_{(k,l)}(z))\\
\iff\\
(\mathcal{G}_1 \coprod \mathcal{G}_2)(\sigma)(\iota^{(k,l)}_1 \circ f_{(k,l)}(z)) \sim_{\star} (\mathcal{G}_1 \coprod \mathcal{G}_2)(\sigma)(\iota^{(k,l)}_2 \circ g_{(k,l)}(z))\\
\iff\\
(\mathcal{G}_1(\sigma) \coprod \mathcal{G}_2(\sigma))(\iota^{(k,l)}_1 \circ f_{(k,l)}(z)) \sim_{\star} (\mathcal{G}_1(\sigma) \coprod \mathcal{G}_2(\sigma))(\iota^{(k,l)}_2 \circ g_{(k,l)}(z))\\
\iff\\
\iota^{\star}_1(\mathcal{G}_1(\sigma)(f_{(k,l)}(z))) \sim_{\star} \iota^{\star}_2(\mathcal{G}_2(\sigma)(g_{(k,l)}(z)))
\end{gather*}
By naturality of $f$ and $g$, we have that $\mathcal{G}_1(\sigma) \circ f_{(k,l)} = f_{\star} \circ \mathcal{G}_0(\sigma)$ and $\mathcal{G}_2(\sigma) \circ g_{(k,l)} = g_{\star} \circ \mathcal{G}_0(\sigma)$, so we have
\begin{gather*}
\iota^{\star}_1(\mathcal{G}_1(\sigma)(f_{(k,l)}(z))) \sim_{\star} \iota^{\star}_2(\mathcal{G}_2(\sigma)(g_{(k,l)}(z)))\\
\iff\\
\iota^{\star}_1(f_{\star}(\mathcal{G}_0(\sigma)(z))) \sim_{\star} \iota^{\star}_2(g_{\star}(\mathcal{G}_0(\sigma)(z))).
\end{gather*}
The last line holds by definition of $\sim_{\star}$.
The rest follows by a more detailed case analysis on $w_1 \sim_{(k,l)} w_2$.
Checking that $\tilde{\alpha}$ preserves identities and composition, so that $\mathcal{G}_1 \coprod_{\mathcal{G}_0} \mathcal{G}_2$ is indeed a functor, is routine.

\section{Frobenius isomorphism}
\label{sec:appendix-frobenius-isomorphism}

\begin{theorem}[Theorem 4.1~\cite{bonchi_string_2022-1}]
  The functor $\llangle - \rrangle : \mathbf{H}_{\Sigma} \to \catname{HypI}(\Sigma)$ is an isomorphism of PROPs, i.e., $\mathbf{H}_{\Sigma} \cong \catname{HypI}(\Sigma)$.
\end{theorem}
\begin{proof}
We present a sketch of the proof here, and refer the reader to~\cite{bonchi_string_2022-1} for the full details.
The idea is to show that $\catname{HypI}(\Sigma)$ has the universal property of the coproduct, namely that for any PROP $\mathbb{A}$ and PROP morphisms $\alpha : \mathbf{S}_{\Sigma} \to \mathbb{A}$ and $\beta : \catname{Frob} \to \mathbb{A}$, there exists a unique PROP morphism $\gamma$ such that the following diagram commutes:
\[\begin{tikzcd}
	{\mathbf{S}_{\Sigma}} && {\catname{HypI}(\Sigma)} && {\catname{Frob}} \\
	\\
	&& {\mathbb{A}}
	\arrow["{\llbracket - \rrbracket}", from=1-1, to=1-3]
	\arrow["\alpha"', from=1-1, to=3-3]
	\arrow["\gamma", from=1-3, to=3-3]
	\arrow["{[-]}"', from=1-5, to=1-3]
	\arrow["\beta"', from=1-5, to=3-3]
\end{tikzcd}~.\]
As all the arrows are identity-on-objects functors, it suffices to show that every morphism in $\catname{HypI}(\Sigma)$ can be decomposed in a unique way into components such that every component is in the image of either $\llbracket - \rrbracket$ or $[-]$.
Let $m \xrightarrow{f} \mathcal{G} \xleftarrow{g} n$ be a morphism in $\catname{HypI}(\Sigma)$ and let $\mathcal{G} = (V_{\mathcal{G}}, E_{\mathcal{G}}, s, t, l)$ be its carrier hypergraph.
Let $e_1, \ldots, e_n$ be some ordering on $E_{\mathcal{G}}$.
Let $\tilde{m} \xrightarrow{\tilde{f}} \tilde{E} \xleftarrow{\tilde{g}} \tilde{n}$ be a cospan with the carrier formed by edges from $\mathcal{G}$, disconnected from each other and stacked vertically with $\tilde{m}$ being the collection of all input vertices for each edge and $\tilde{m}$ being the collection of all output vertices for each edge.
There are functions $i : n \to V_{\mathcal{G}}$,  $\tilde{i} : \tilde{n} \to V_{\mathcal{G}}$, $j : m \to V_{\mathcal{G}}$, and $\tilde{j} : \tilde{m} \to V_{\mathcal{G}}$ that map vertices in the interfaces to their occurrences in the set of all vertices.
The decomposition of $m \xrightarrow{f} \mathcal{G} \xleftarrow{g} n$ is then given by the composition of cospans
% \begin{gather}\label{eq:decomposition}
% \displaystyle (m \xrightarrow{i} V_{\mathcal{G}} \xleftarrow{[\id, \tilde{j}]} V_{\mathcal{G}} \otimes \tilde{m}) \fatsemi (V_{\mathcal{G}} \otimes \tilde{m} \xrightarrow{\id \otimes \tilde{f}} V_{\mathcal{G}} \otimes \tilde{E} \xleftarrow{\id \otimes \tilde{g}} V_{\mathcal{G}} \otimes \tilde{n}) \fatsemi (V_{\mathcal{G}} \otimes \tilde{n} \xrightarrow{[\id, \tilde{i}]} V_{\mathcal{G}} \xleftarrow{i} n)~.
% \end{gather}
\begin{equation}\label{eq:decomposition}
  \begin{adjustbox}{scale=0.9}
\begin{tikzcd}
	m & {V_{\mathcal{G}}} & {V_{\mathcal{G}}\otimes\tilde{m}} & {V_{\mathcal{G}}\otimes\tilde{E}} & {V_{\mathcal{G}}\otimes\tilde{n}} & {V_{\mathcal{G}}} & n
	\arrow["i", from=1-1, to=1-2]
	\arrow["{{[\id,\tilde{j}]}}"', from=1-3, to=1-2]
	\arrow["{{\id\otimes\tilde{f}}}", from=1-3, to=1-4]
	\arrow["{{\id\otimes\tilde{g}}}"', from=1-5, to=1-4]
	\arrow["{{[\id,\tilde{i}]}}", from=1-5, to=1-6]
	\arrow[from=1-6, to=1-5]
	\arrow["i"', from=1-7, to=1-6]
\end{tikzcd}
\end{adjustbox}
\end{equation}
For example, below is a cospan (on the left) and its decomposition (on the right); we show the embedding of the interfaces using dashed arrows to avoid excessive cluttering of the cospans and use the labels to assign identities to vertices in $V_{\mathcal{G}}$.
\[
\adjustbox{scale=0.8}{\tikzfig{./figures/string/piling_example_1}} \quad \cong \adjustbox{scale=0.8}{\tikzfig{./figures/string/piling_example_3}} \fatsemi\; \adjustbox{scale=0.8}{\tikzfig{./figures/string/piling_example_2}} \fatsemi \adjustbox{scale=0.8}{\tikzfig{./figures/string/piling_example_4}}~.
\]
By computing the pushouts for the compositions on the right, we can verify that the composition is indeed isomorphic to the cospan on the left. 
Intuitively, the middle cospan records all edges and vertices of the carrier hypergraph and the cospans on the sides record the incidence of the hyperedges and vertices.
For example, in the first composition, vertices $u_1, u_3$ of hyperedges $a$ and $b$ get identified with the vertex $z_1$ according to the interface embedding.

The middle cospan of the decomposition~\ref{eq:decomposition} is in the image of $\llbracket - \rrbracket$ as it is a tensor of identities and generators from $\Sigma$.
The cospans on the sides are essentially cospans of finite sets and hence in the image of $[-]$ as $\catname{Frob} \cong \catname{Cospan}(\mathbb{F})$---a result of Lack~\cite{lackDedicatedAurelioCarboni}.
This means that every morphism $\gamma : \catname{HypI}(\Sigma) \to \mathbb{A}$ is uniquely determined by its action on the components of the decomposition of each morphism.
\end{proof}

\section{Pushout complements}
\label{sec:appendix-pushout-complement}

\begin{proposition}
  Let $K \xrightarrow{i} L \xrightarrow{f} G$ be morphisms in $\catname{Hyp}(\Sigma)$.
  A pushout complement of $K \xrightarrow{i} L \xrightarrow{f} G$ exists if and only if the no-dangling and no-identification conditions are satisfied.
\end{proposition}
\begin{proof}
  For the if direction, suppose the two conditions hold, and define $L^{\bot}$ as the following sub-hypergraph of $G$
  \[
  L^{\bot} \defeq (V_{L^{\bot}}, E_{L^{\bot}}, s_{L^{\bot}}, t_{L^{\bot}}, l_{L^{\bot}}) \hookrightarrow G, 
  \]
  where 
  \begin{gather*}
    V_{L^{\bot}} \defeq \bigl(V_G \setminus f(V_L)\bigr) \cup f(i(V_K))\\
    E_{L^{\bot}} \defeq \bigl(E_G \setminus f(E_L)\bigr) \cup f(i(E_K))\\
    s_{L^{\bot}} \defeq s_G \restriction_{E_{L^{\bot}}} \quad t_{L^{\bot}} \defeq t_G \restriction_{E_{L^{\bot}}} \quad l_{L^{\bot}} \defeq l_G \restriction_{E_{L^{\bot}}}
  \end{gather*}
  This is a well-defined sub-hypergraph: a hyperedge of $L^{\bot}$ either lies outside the image of $f$, in which case the no-dangling condition places each of its sources and targets outside the image of $f$ or inside the image of $f \circ i$, hence in $V_{L^{\bot}}$; or it lies in the image of $f \circ i$, in which case its sources and targets lie in the image of $f \circ i$ as well, since $f \circ i$ is a homomorphism.
  Let $g \colon L^{\bot} \hookrightarrow G$ be the inclusion and $c \colon K \to L^{\bot}$ the corestriction of $f \circ i$, which lands in $L^{\bot}$ by construction; then $g \circ c = f \circ i$, so the square commutes.

  It remains to show the square is a pushout.
  Let $P \defeq L \coprod_{K} L^{\bot}$ be the pushout of $L \xleftarrow{i} K \xrightarrow{c} L^{\bot}$ computed as in~\cref{def:pushout_hyp}: the vertices and edges of the resulting hypergraph are quotients $V_{L} \coprod_{V_{K}} V_{L^{\bot}} / \sim$ and $E_{L}\coprod_{E_{K}} E_{L^{\bot}} / \sim$ by the least equivalence relation identifying $\iota_{1}^{V}(i_{V}(z))$ with $\iota_{2}^{V}(c_{V}(z))$ for $z \in V_{K}$ and similarly for edges.
  Let $u \colon P \to G$ be the homomorphism induced by $f$ and $g$ via the universal property, given explicitly on vertices by $u_{V}([\iota^{V}_{1}(w)]) = f_{V}(w)$ and $u_{V}([\iota^{V}_{2}(z)]) = g_{V}(z)$, and analogously on edges.
  We show that $u_{V}$ and $u_{E}$ are bijections; since a homomorphism that is bijective on vertices and on edges is an isomorphism of hypergraphs, this suffices.

  For surjectivity, a vertex of $G$ either lies in $f(V_{L})$, hence in the image of $u_{V} \circ [\iota^{V}_{1}]$, or lies in $V_{G} \setminus f(V_{L}) \subseteq V_{L^{\bot}}$, hence in the image of $u_{V} \circ [\iota^{V}_{2}]$; the same argument applies to edges.

  For injectivity, note that $u_{V}$ is injective if and only if any two elements of $V_{L} \coprod V_{L^{\bot}}$ with the same image under the copairing $[f_{V}, g_{V}] \colon V_{L} \coprod V_{L^{\bot}} \to V_{G}$ are related by $\sim$: the copairing factors as $u_{V}$ after the quotient projection.
  So suppose two elements of $V_{L} \coprod V_{L^{\bot}}$ have the same image under $[f_{V}, g_{V}]$; we show that they are related by $\sim$.
  There are three cases.
  If both are of the form $\iota^{V}_{2}(z_{1}), \iota^{V}_{2}(z_{2})$, then $g_{V}(z_{1}) = g_{V}(z_{2})$ implies $z_{1} = z_{2}$, since $g$ is an inclusion.
  If both are of the form $\iota^{V}_{1}(w_{1}), \iota^{V}_{1}(w_{2})$ with $w_{1} \neq w_{2}$, then $f_{V}(w_{1}) = f_{V}(w_{2})$, so by the no-identification condition $w_{1} = i_{V}(z_{1})$ and $w_{2} = i_{V}(z_{2})$ for some $z_{1}, z_{2} \in V_{K}$.
  Then $c_{V}(z_{1}) = f_{V}(i_{V}(z_{1})) = f_{V}(w_{1}) = f_{V}(w_{2}) = c_{V}(z_{2})$, where the outer equalities hold because $c$ is the corestriction of $f \circ i$, and hence
  \[
  \iota^{V}_{1}(w_{1}) = \iota^{V}_{1}(i_{V}(z_{1})) \sim \iota^{V}_{2}(c_{V}(z_{1})) = \iota^{V}_{2}(c_{V}(z_{2})) \sim \iota^{V}_{1}(i_{V}(z_{2})) = \iota^{V}_{1}(w_{2}).
  \]
 Finally, suppose the elements are $\iota^{V}_{1}(w)$ and $\iota^{V}_{2}(z)$ with $f_{V}(w) = g_{V}(z)$.
  Since $g_{V}(z) \in V_{L^{\bot}}$ and $f_{V}(w) \in f(V_{L})$, the construction of $V_{L^{\bot}}$ yields $f_{V}(w) \in f(i(V_{K}))$, say $f_{V}(w) = f_{V}(i_{V}(y))$ for some $y \in V_{K}$.
  If $w = i_{V}(y)$, then
  \[
  g_{V}(c_{V}(y)) = f_{V}(i_{V}(y)) = f_{V}(w) = g_{V}(z),
  \]
  where the first equality is the commutativity $g \circ c = f \circ i$; injectivity of $g_{V}$ then gives $c_{V}(y) = z$, and hence $\iota^{V}_{1}(w) = \iota^{V}_{1}(i_{V}(y)) \sim \iota^{V}_{2}(c_{V}(y)) = \iota^{V}_{2}(z)$.
  Otherwise, $w$ and $i_{V}(y)$ are distinct vertices identified by $f$, so by the no-identification condition $w = i_{V}(y')$ for some $y' \in V_{K}$, and the previous argument applies with $y'$ in place of $y$.
  In all cases the two elements are related by $\sim$, so $u_{V}$ is injective; the argument for $u_{E}$ is identical, using the no-identification condition for edges in the second case.
  Hence $u$ is an isomorphism, the constructed square is a pushout, and $L^{\bot}$ is a pushout complement of $K \xrightarrow{i} L \xrightarrow{f} G$.

  For the converse, suppose that a pushout complement $K \xrightarrow{c} L^{\bot} \xrightarrow{g} G$ exists but at least one of the conditions fails.
  Since pushouts are unique up to an isomorphism commuting with the pushout injections, we may by~\cref{def:pushout_hyp} assume that $V_{G} = V_{L} \coprod V_{L^{\bot}} / {\sim}$ and $E_{G} = E_{L} \coprod E_{L^{\bot}} / {\sim}$, with $f$ and $g$ the composites of the coproduct injections with the quotient projections.
  We first record a consequence of this description: if $w \in V_{L}$ and $\iota^{V}_{1}(w)$ is related by $\sim$ to any element of the coproduct other than itself, then $w \in i(V_{K})$.
  Indeed, every generating pair of $\sim$ is of the form $\iota^{V}_{1}(i_{V}(z)) \sim \iota^{V}_{2}(c_{V}(z))$, so a non-trivial chain of generating pairs passing through $\iota^{V}_{1}(w)$ must contain a generator whose left component is $\iota^{V}_{1}(w)$ itself, forcing $w = i_{V}(z)$ for some $z \in V_{K}$.
  The same holds for edges.

  Suppose the no-identification condition fails: there are distinct items of $L$, say vertices $w_{1} \neq w_{2}$ (the case of edges is identical), with $f_{V}(w_{1}) = f_{V}(w_{2})$ and at least one of them outside the image of $i$.
  Then $\iota^{V}_{1}(w_{1})$ and $\iota^{V}_{1}(w_{2})$ are distinct elements of the coproduct in the same $\sim$-class, so both are related by $\sim$ to an element other than themselves, and by the consequence above both lie in the image of $i_{V}$---a contradiction.

  Suppose the no-dangling condition fails: there is a hyperedge $e \in E_{G}$ outside the image of $f$ with a source or target $v \in f(V_{L}) \setminus f(i(V_{K}))$, say $v = f_{V}(w)$ with $w \notin i(V_{K})$.
  Since $f$ and $g$ are jointly surjective on edges and $e \notin f(E_{L})$, we have $e = g_{E}(e')$ for some $e' \in E_{L^{\bot}}$.
  As $g$ is a homomorphism, $g_{V}^{*}(s_{L^{\bot}}(e')) = s_{G}(e)$ and likewise for targets, so the corresponding source or target $z$ of $e'$ satisfies $g_{V}(z) = v = f_{V}(w)$.
  Then $\iota^{V}_{1}(w)$ and $\iota^{V}_{2}(z)$ are distinct elements of the coproduct in the same $\sim$-class, so $\iota^{V}_{1}(w)$ is related by $\sim$ to an element other than itself, and by the consequence above $w \in i(V_{K})$, so we, again, reach a contradiction.
\end{proof}


\section{DPOI rewriting}
\label{sec:appendix-dpoi-rewriting-frob}

\begin{theorem}[{\cite[Thm. 4.9]{bonchi_string_2022-1}}]
  Let $\langle l, r \rangle$ be a rewrite rule in $\mathbf{H}_{\Sigma}$.
  Then,
  \[
  a \rewriteArrow{\langle l, r \rangle} b \quad \iff \quad \llangle \corner{a} \rrangle \rightsquigarrow_{\langle \llangle \corner{l} \rrangle, \llangle \corner{r} \rrangle \rangle} \llangle \corner{b} \rrangle
  \]
\end{theorem}
\begin{proof}
For the only-if direction, $a \rewriteArrow{\langle l, r \rangle} b$ by definition means that there exists a decomposition of $a$ and $b$ as
\[
\adjustbox{scale=0.7}{\tikzfig{./figures/string/decomposition}}~. 
\]
Applying $\corner{-}$ to the above gives us the decomposition of $\corner{a}$ as
\[
\adjustbox{scale=0.7}{\tikzfig{./figures/string/cup_decomposition}}~,
\]
or, by abbreviating
\[
\adjustbox{scale=0.7}{\tikzfig{./figures/string/cup_decomposition_2}}~,
\]
as
\[
\adjustbox{scale=0.7}{\tikzfig{./figures/string/cup_decomposition_3}}~.
\]
Analogously, we also get
\[
\adjustbox{scale=0.7}{\tikzfig{./figures/string/cup_decomposition_4}}~.
\]
The dashed lines result in the following
\begin{equation}
\corner{a} = \corner{l} \fatsemi \tilde{a} \qquad \corner{b} = \corner{r} \fatsemi \tilde{a}\label{eq:tilde_decomposition}~,
\end{equation}
where $\tilde{a} = (\id_{i} \otimes \cup_{k} \otimes \id_{j}) \fatsemi (a_1^{\star} \otimes a_2)$.
Applying the functor $\llangle - \rrangle$ to the above gives us the following cospans
\begin{gather*}\label{gather:cospans_frobenius_rewriting}
\llangle \corner{l} \rrangle = 0 \xrightarrow{} L \xleftarrow{} i \coprod j \qquad \llangle \corner{r} \rrangle = 0 \xrightarrow{} R \xleftarrow{} i \coprod j\\
\llangle \tilde{a} \rrangle = i \coprod j \xrightarrow{} L^{\bot} \xleftarrow{} m \coprod n\\
\llangle \corner{a} \rrangle = 0 \xrightarrow{} G \xleftarrow{} m \coprod n \qquad \llangle \corner{b} \rrangle = 0 \xrightarrow{} H \xleftarrow{} m \coprod n~.
\end{gather*}
By functoriality and by~\cref{eq:tilde_decomposition}, we have
\[
\llangle \corner{a} \rrangle = \llangle \corner{l} \rrangle \fatsemi \llangle \tilde{a} \rrangle \qquad \llangle \corner{b} \rrangle = \llangle \corner{r} \rrangle \fatsemi \llangle \tilde{a} \rrangle~,
\]
since composition of cospans is given by pushouts, we get the diagram
\[\begin{tikzcd}
	L && {i \coprod j} && R \\
	\\
	G && {L^{\bot}} && H \\
	&& {m \coprod n}
	\arrow[from=1-1, to=3-1]
	\arrow[from=1-3, to=1-1]
	\arrow[from=1-3, to=1-5]
	\arrow[from=1-3, to=3-3]
	\arrow[from=1-5, to=3-5]
	\arrow["\lrcorner"{anchor=center, pos=0.125, rotate=90}, draw=none, from=3-1, to=1-3]
	\arrow[from=3-3, to=3-1]
	\arrow[from=3-3, to=3-5]
	\arrow["\lrcorner"{anchor=center, pos=0.125, rotate=180}, draw=none, from=3-5, to=1-3]
	\arrow[from=4-3, to=3-1]
	\arrow[from=4-3, to=3-3]
	\arrow[from=4-3, to=3-5]
\end{tikzcd}\]
and by~\cref{def:dpoi_rewriting} we have $\llangle \corner{a} \rrangle \rightsquigarrow_{\langle \llangle \corner{l} \rrangle, \llangle \corner{r} \rrangle \rangle} \llangle \corner{b} \rrangle$.

For the if direction, suppose $\llangle \corner{a} \rrangle \rightsquigarrow_{\langle \llangle \corner{l} \rrangle, \llangle \corner{r} \rrangle \rangle} \llangle \corner{b} \rrangle$.
Then, by~\cref{def:dpoi_rewriting}, there exists a pushout complement $L^{\bot}$ such that the diagram as above commutes with two squares being pushouts.
We denote the cospans involved in the DPOI square as in~\cref{gather:cospans_frobenius_rewriting}.
Given $i \coprod j \xrightarrow{} L^{\bot} \xleftarrow{} m \coprod n$ pick $\hat{a}$ in $\mathbf{H}_{\Sigma}$ such that $\llangle \hat{a} \rrangle = i \coprod j \xrightarrow{} L^{\bot} \xleftarrow{} m \coprod n$, which exists by fullness of $\llangle - \rrangle$.
Because composition in $\catname{HypI}(\Sigma)$ is given by pushouts, we have
\begin{equation}\label{eq:only_if_direction}
\adjustbox{scale=0.85}{$\displaystyle
\begin{gathered}
  \llangle \corner{a} \rrangle = (0 \to G \leftarrow m \coprod n)
  = (0 \to L \leftarrow i \coprod j) \fatsemi (i \coprod j \to L^{\bot} \leftarrow m \coprod n)
  = \llangle \corner{l} \rrangle \fatsemi \llangle \hat{a} \rrangle = \llangle \corner{l} \fatsemi \hat{a} \rrangle\\
  \llangle \corner{b} \rrangle = (0 \to H \leftarrow m \coprod n)
  = (0 \to R \leftarrow i \coprod j) \fatsemi (i \coprod j \to L^{\bot} \leftarrow m \coprod n)
  = \llangle \corner{r} \rrangle \fatsemi \llangle \hat{a} \rrangle = \llangle \corner{r} \fatsemi \hat{a} \rrangle
\end{gathered}
$}
\end{equation}
By faithfulness of $\llangle - \rrangle$, the decompositions in~\cref{eq:only_if_direction} lift to decompositions of $\corner{a} = \corner{l} \fatsemi \hat{a}$ and $\corner{b} = \corner{r} \fatsemi \hat{a}$:
\[
\adjustbox{scale=0.8}{\tikzfig{./figures/string/cup_decomposition_5}}~.
\]
By inverting $\corner{-}$, we get the decompositions of $a$ and $b$ as
\[
\adjustbox{scale=0.8}{\tikzfig{./figures/string/cup_decomposition_6}}~,
\]
with dashed lines outlining the parts relevant to~\cref{def:rewriting_frob}, i.e. the part above the top dashed line corresponds to $a_2$ and the part below the bottom dashed line corresponds to $a_1$. 
This, by definition, means that $a \rewriteArrow{\langle l, r \rangle} b$.
\end{proof}


\section{Convex DPOI rewriting}
\label{sec:appendix-convex-dpoi-rewriting}

\begin{theorem}[{\cite[Thm. 35]{bonchi_string_2022-2}}]
Let $\mathcal{R}$ be any rewriting system on $\mathbf{S}_{\Sigma}$, then
\[
a \Rightarrow_{\mathcal{R}} b \qquad \iff \qquad \llangle \corner{a} \rrangle \Rrightarrow_{\llangle \corner{\mathcal{R}} \rrangle} \llangle \corner{b} \rrangle~.
\]
\end{theorem}
\begin{proof}
For the only-if direction,~\cref{theorem:dpoi_rewriting_frob} implies that 
\[
\llangle \corner{a} \rrangle \rightsquigarrow_{\llangle \corner{\mathcal{R}} \rrangle} \llangle \corner{b} \rrangle~,
\]
which gives us an analogous decomposition as in the proof of~\cref{theorem:dpoi_rewriting_frob}.
In particular, we have that
\[
\big\llangle \adjustbox{scale=0.7}{\tikzfig{./figures/string/convex_cup_decomposition}}\big\rrangle = i \coprod j \xrightarrow{} L^{\bot} \xleftarrow{} m \coprod n~.
\]
We can check that $m \coprod j \xrightarrow{} L^{\bot} \xleftarrow{} n \coprod i$ is an mda-cospan by \enquote{unbending} the wires:
\[
\adjustbox{scale=0.7}{\tikzfig{./figures/string/convex_cup_decomposition_2}}~.
\]
This is indeed monogamous as $a_1$ and $a_2$ are morphisms in $\mathbf{S}_{\Sigma}$.
Then, we have
\[
\llangle \corner{a} \rrangle
= (0 \xrightarrow{} G \xleftarrow{} m \coprod n)
= (0 \xrightarrow{} L \xleftarrow{} i \coprod j) \fatsemi (i \coprod j \xrightarrow{} L^{\bot} \xleftarrow{} m \coprod n)
= \llangle \corner{l} \rrangle \fatsemi \llangle \tilde{a} \rrangle
\]
and let $f : L \to G$ be the pushout injection of the composition; we check that $f$ is mono and that its image is a convex sub-hypergraph of $G$.

For monicity, recall that the composition is computed by the pushout of
$L \xleftarrow{[a_1,a_2]} i \coprod j \xrightarrow{[c_1,c_2]} L^{\bot}$,
of which $f$ is the injection opposite to $[c_1,c_2]$.
Since $[c_1,c_2]$ is mono, as established when verifying the boundary
complement conditions, and pushouts along monos in $\catname{Hyp}(\Sigma)$
are mono, $f$ is mono.

For convexity, first observe that, by monogamy of $G$, a path that leaves
$f(L)$ can only do so through a vertex shared with $L^{\bot}$ whose unique
out-hyperedge lies in $L^{\bot}$, and can only re-enter $f(L)$ through a
shared vertex whose unique in-hyperedge lies in $L^{\bot}$.
The shared vertices are exactly $i \coprod j$; every $i$-vertex has its
out-hyperedge in $f(L)$, and every $j$-vertex has its in-hyperedge in
$f(L)$, so convexity reduces to the claim that $L^{\bot}$ contains no path
from a $j$-vertex to an $i$-vertex.
This is where we use that $\llangle \tilde{a} \rrangle$ arises from the
decomposition: the carrier $L^{\bot}$ is the carrier of $\llangle a_1
\rrangle$ glued to the carrier of $\llangle a_2 \rrangle$ along the
$k$-vertices only, with the $i$-vertices lying in the $a_1$-part and the
$j$-vertices in the $a_2$-part.
Assign rank $0$ to the hyperedges of the $a_1$-part and rank $1$ to those
of the $a_2$-part.
By monogamy of the two layers, every vertex of $L^{\bot}$ whose in- and
out-hyperedges both exist within $L^{\bot}$ has the in-hyperedge of rank
at most the out-hyperedge: internal vertices of a layer have both in the
same layer, and $k$-vertices have their in-hyperedge in the $a_1$-part and
out-hyperedge in the $a_2$-part.
Hence ranks are non-decreasing along paths.
A path starting at a $j$-vertex begins with a hyperedge of rank $1$, since
$j$-vertices are inputs of the $a_2$-part and carry no hyperedge of the
$a_1$-part; a path ending at an $i$-vertex ends with a hyperedge of rank
$0$, since the unique in-hyperedge of an $i$-vertex within $L^{\bot}$ lies
in the $a_1$-part.
No path can therefore lead from a $j$-vertex to an $i$-vertex, and $f(L)$
is convex.
That is, we have that $\llangle \corner{a} \rrangle \Rrightarrow_{\llangle \corner{\mathcal{R}} \rrangle} \llangle \corner{b} \rrangle$.

For the if direction, suppose $\llangle \corner{a} \rrangle \Rrightarrow_{\llangle \corner{\mathcal{R}} \rrangle} \llangle \corner{b} \rrangle$.
Then, by definition, there exists a convex match $f : L \to G$ and a boundary pushout complement $i \coprod j \xrightarrow{} L^{\bot} \xleftarrow{} m \coprod n$ such that the diagram commutes and the marked squares are pushouts, for some rule $\langle \llangle \corner{l} \rrangle, \llangle \corner{r} \rrangle \rangle$ in $\llangle \corner{\mathcal{R}} \rrangle$ with $\langle l, r \rangle : i \to j$ in $\mathcal{R}$.

Since $a$ is a morphism of $\mathbf{S}_{\Sigma}$, the unfolded cospan $m \xrightarrow{q_1} G \xleftarrow{q_2} n = \llbracket a \rrbracket$ is an mda-cospan by~\cref{theorem:monogamous_isomorphism}.
The match $f$ is mono with convex image, so~\cref{lemma:convex_decomposition} yields $k \in \mathbb{N}$ and a factorisation into mda-cospans
\[
(m \xrightarrow{} C_1 \xleftarrow{} k \coprod i)
\;\fatsemi\;
\bigl((k \xrightarrow{\id} k \xleftarrow{\id} k) \oplus (i \xrightarrow{} L \xleftarrow{} j)\bigr)
\;\fatsemi\;
(k \coprod j \xrightarrow{} C_2 \xleftarrow{} n)~,
\]
where the middle factor restricts to the unique cospan structure on $f(L) \cong L$ guaranteed by the lemma, compatible with the interface $i \coprod j \to L$ of the rule.
Applying~\cref{theorem:monogamous_isomorphism} once more, this time using fullness, we obtain morphisms $a_1 : m \to k \coprod i$ and $a_2 : k \coprod j \to n$ in $\mathbf{S}_{\Sigma}$ with
\[
\llbracket a_1 \rrbracket = (m \xrightarrow{} C_1 \xleftarrow{} k \coprod i)
\qquad
\llbracket a_2 \rrbracket = (k \coprod j \xrightarrow{} C_2 \xleftarrow{} n)~.
\]
By functoriality, $\llbracket a \rrbracket = \llbracket a_1 \fatsemi (\id_k \otimes l) \fatsemi a_2 \rrbracket$, and since $\llbracket - \rrbracket$ is faithful, $a = a_1 \fatsemi (\id_k \otimes l) \fatsemi a_2$ already in $\mathbf{S}_{\Sigma}$.
This exhibits a redex, so the rule applies syntactically: setting $b' \defeq a_1 \fatsemi (\id_k \otimes r) \fatsemi a_2$, we have $a \Rightarrow_{\mathcal{R}} b'$.

It remains to show that $b' = b$.
The factorisation above assembles, exactly as in the only-if direction, into a second convex DPOI square for the same match $f$ and the same rule leg $i \coprod j \to L$, whose pushout complement is the folded context $i \coprod j \xrightarrow{} L^{\bot\prime} \xleftarrow{} m \coprod n$ built from $a_1$ and $a_2$, and whose result is $\llangle \corner{b'} \rrangle$.
By~\cref{prop:boundary_pushout_complement_unique}, boundary complements are unique, so $L^{\bot\prime}$ and $L^{\bot}$ are isomorphic as pushout complements, i.e.\ via an isomorphism commuting with the arrows from $i \coprod j$ and to $G$.
Pushing out $R \xleftarrow{} i \coprod j$ along the two isomorphic complements yields isomorphic cospans, hence
\[
\llangle \corner{b} \rrangle = (0 \xrightarrow{} E \xleftarrow{} m \coprod n) = \llangle \corner{b'} \rrangle~.
\]
Since $\llangle - \rrangle$ is an isomorphism of PROPs by~\cref{theorem:frobenius_isomorphism}, it follows that $\corner{b} = \corner{b'}$ in $\mathbf{H}_{\Sigma}$; unfolding, and using that $\mathbf{S}_{\Sigma} \to \mathbf{H}_{\Sigma}$ is faithful, $b = b'$ in $\mathbf{S}_{\Sigma}$.
Hence $a \Rightarrow_{\mathcal{R}} b$, as required.
\end{proof}

\chapter{Additional material for Chapter~\ref{chapter:monoidal}}
\label{sec:appendix-monoidal-figures}

\begin{figure}
\[
\adjustbox{max width=\textwidth}{
    \tikzfig{./figures/appendix/egraph-translation-step-by-step-b-c}
}
\]
\caption{Transformation of string diagram (b) into string diagram (c) via the rewrite rule $(x * y) / z \to x * (y / z)$ from~\cref{fig:egg_egraph_string}.}
\label{fig:appendix-egraph-translation-step-by-step-b-c}
\end{figure}

\Cref{fig:appendix-egraph-translation-step-by-step-b-c} provides a futher elaboration on the transformation of string diagram (b) into string diagram (c) in~\cref{fig:egg_egraph_string} via the rewrite rule $(x * y) / z \to x * (y / z)$.
The first row shows how $\catname{SLat}$-equations are applied to reveal the redex for the rule.
Then the content of the dashed box is appropriately unified using the naturality of $\counit$ so that common structure can be factored out.
Since in this setting morphisms have no variables, the unification is always possible.
The last row futher compacts the diagram to match the corresponding e-graph (c) in~\cref{fig:egg_egraph}.

\section{Pushout complements}
\label{sec:appendix-monoidal-pushout-complements}

\begin{lemma}[Shape of a complement]
Let $i \coprod j \xrightarrow{f} \mathcal{L} \xrightarrow{m} \mathcal{G}$ be as
in~\cref{prop:e_pushout_complement_exists} and let
$i \coprod j \xrightarrow{c} \mathcal{L}^{\bot} \xrightarrow{l^{\bot}} \mathcal{G}$
be a pushout complement.
Then
\begin{enumerate}
  \item $l^{\bot}$ is strict;
  \item $l^{\bot}_{V}\bigl(V_{\mathcal{L}^{\bot}}\bigr)
        = \bigl(V_{\mathcal{G}} \setminus m(V_{\mathcal{L}})\bigr) \cup m\bigl(f(V_{i \coprod j})\bigr)$
        and
        $l^{\bot}_{E}\bigl(E_{\mathcal{L}^{\bot}}\bigr)
        = \bigl(E_{\mathcal{G}} \setminus m(E_{\mathcal{L}})\bigr) \cup \{a\}$;
  \item if $f$ is a monomorphism, then so is $l^{\bot}$.
\end{enumerate}
\end{lemma}
\begin{proof}
Throughout we may assume, by the observation made in the proof of~\cref{prop:pushout_complement_exists}, that $\mathcal{G}$ is the constructed pushout, so
that $V_{\mathcal{G}} = V_{\mathcal{L}} \coprod V_{\mathcal{L}^{\bot}} / R_{V}$
and $E_{\mathcal{G}} = E_{\mathcal{L}} \coprod E_{\mathcal{L}^{\bot}} / R_{E}$,
with $m$ and $l^{\bot}$ the composites of the coproduct injections with the
quotient projections, and that every generating pair of $R_{V}$, respectively
$R_{E}$, has the form
$\bigl(\iota_{1}(f(z)), \iota_{2}(c(z))\bigr)$.

\begin{enumerate}
\item The span
$\mathcal{L} \xleftarrow{f} i \coprod j \xrightarrow{c} \mathcal{L}^{\bot}$
has $f$ strict and $c$ anchored, so by~\cref{prop:pushout_exists} the
coprojection out of $\mathcal{L}^{\bot}$ is the strict one.

\item For the inclusion $\supseteq$: an element outside the image of $m$ lies in
the image of $l^{\bot}$ by joint surjectivity; and
$f \fatsemi m = c \fatsemi l^{\bot}$ puts $m(f(V_{i \coprod j}))$ and
$a = m_{E}(f_{E}(\top_{i \coprod j}))$ in the image of $l^{\bot}$.
For $\subseteq$, consider $y = l^{\bot}(x)$, if $y$ is a vertex and $y \not \in m(V_{\mathcal{L}})$ then $y \in V_{\mathcal{G}} \setminus m(V_{\mathcal{L}})$ and similarly if $y$ is an edge. 
So assume $y = m(w)$ for $w$ in $\mathcal{L}$ and suppose $x$ and $w$ are vertices.
Since they are identified in the quotient, there must be an element $z \in V_{i \coprod j}$ such that $\iota_{1}(f(z)) = x$ and $\iota_{2}(c(z)) = w$ and therefore $l^{\bot}(x) \in m(f(V_{i \coprod j}))$.
Now if $y$ was an edge, then $y = a$.

\item Since $f$ is monic, its vertex and edge components are injective. Fix one
sort $S \in \{V,E\}$ and suppose
$l^{\bot}_{S}(x_{1}) = l^{\bot}_{S}(x_{2})$. Then
$\iota_{2}(x_{1})$ and $\iota_{2}(x_{2})$ are joined by an alternating
$R_{S}$-chain. Every portion of this chain passing from the
$\mathcal{L}^{\bot}$-summand back to that summand has the form
\[
\iota_{2}(c_{S}(z_{1}))
\sim
\iota_{1}(f_{S}(z_{1}))
=
\iota_{1}(f_{S}(z_{2}))
\sim
\iota_{2}(c_{S}(z_{2})).
\]
Injectivity of $f_{S}$ gives $z_{1}=z_{2}$, and hence
$c_{S}(z_{1})=c_{S}(z_{2})$. Thus every term of the chain lying in the
$\mathcal{L}^{\bot}$-summand is equal to $x_{1}$, so $x_{2}=x_{1}$.
Consequently both components of $l^{\bot}$ are injective, and therefore
$l^{\bot}$ is a monomorphism.
\end{enumerate}
\end{proof}

\begin{lemma}[Uniqueness of pushout complements]
Let $i \coprod j \xrightarrow{f} \mathcal{L} \xrightarrow{m} \mathcal{G}$ be as
in~\cref{prop:e_pushout_complement_exists} with $f$ a monomorphism, and let
$(\mathcal{L}^{\bot}, c, l^{\bot})$ and $(\mathcal{K}, c', k)$ be pushout
complements of $(f, m)$.
Then there is a unique isomorphism $\phi : \mathcal{L}^{\bot} \to \mathcal{K}$
in $\catname{EHyp}_{\mathrm{s}}(\Sigma)$ with $\phi \fatsemi k = l^{\bot}$, and
it satisfies $c \fatsemi \phi = c'$.
\end{lemma}
\begin{proof}
By~\cref{lemma:complement_shape}, $l^{\bot}$ and $k$ are strict monomorphisms
with the same image, which is described in part (2) of that lemma without
reference to the complement.
Define
\[
\phi_{V} \defeq l^{\bot}_{V} \fatsemi (k_{V})^{-1}~,
\qquad
\phi_{E} \defeq l^{\bot}_{E} \fatsemi (k_{E})^{-1}~,
\]
the inverses being taken on that common image; these are bijections and satisfy
$\phi \fatsemi k = l^{\bot}$ by construction.

\emph{$\phi$ is a homomorphism.}
For sources, using that $l^{\bot}$ and $k$ are homomorphisms,
\[
k_{V}^{*}\bigl(s_{\mathcal{K}}(\phi_{E}(e))\bigr)
= s_{\mathcal{G}}\bigl(k_{E}(\phi_{E}(e))\bigr)
= s_{\mathcal{G}}\bigl(l^{\bot}_{E}(e)\bigr)
= (l^{\bot}_{V})^{*}\bigl(s_{\mathcal{L}^{\bot}}(e)\bigr)
= k_{V}^{*}\bigl(\phi_{V}^{*}(s_{\mathcal{L}^{\bot}}(e))\bigr)~,
\]
and $k_{V}$ is injective, so
$s_{\mathcal{K}}(\phi_{E}(e)) = \phi_{V}^{*}(s_{\mathcal{L}^{\bot}}(e))$.
Targets are identical, and labels follow from
$l_{\mathcal{K}}(\phi_{E}(e)) = l_{\mathcal{G}}(k_{E}(\phi_{E}(e)))
= l_{\mathcal{G}}(l^{\bot}_{E}(e)) = l_{\mathcal{L}^{\bot}}(e)$.

\emph{$\phi$ is strict.}
Both $l^{\bot}$ and $k$ are strict, so
$k_{E}(\phi_{E}(\top_{\mathcal{L}^{\bot}})) = l^{\bot}_{E}(\top_{\mathcal{L}^{\bot}})
= \top_{\mathcal{G}} = k_{E}(\top_{\mathcal{K}})$, and $k_{E}$ is injective,
therefore $\phi_{E}(\top_{\mathcal{L}^{\bot}}) = \top_{\mathcal{K}}$.
For $x \neq \kappa^{E}(\top_{\mathcal{L}^{\bot}})$,
\[
k_{E}\Bigl(\phi_{E}\bigl({<^{\mu}_{\mathcal{L}^{\bot}}}(x)\bigr)\Bigr)
= l^{\bot}_{E}\bigl({<^{\mu}_{\mathcal{L}^{\bot}}}(x)\bigr)
= {<^{\mu}_{\mathcal{G}}}\bigl(l^{\bot}(x)\bigr)
= {<^{\mu}_{\mathcal{G}}}\bigl(k(\phi(x))\bigr)
= k_{E}\bigl({<^{\mu}_{\mathcal{K}}}(\phi(x))\bigr)~,
\]
using strictness of $l^{\bot}$ and of $k$, and injectivity of $k_{E}$ gives
$\phi_{E}({<^{\mu}_{\mathcal{L}^{\bot}}}(x)) = {<^{\mu}_{\mathcal{K}}}(\phi(x))$.
The same argument with the roles of $\mathcal{L}^{\bot}$ and $\mathcal{K}$
exchanged shows that $\phi^{-1} = k \fatsemi (l^{\bot})^{-1}$ is a strict
homomorphism, so $\phi$ is an isomorphism of
$\catname{EHyp}_{\mathrm{s}}(\Sigma)$.

For $z \in V_{i \coprod j} \coprod E_{i \coprod j}$,
\[
k\bigl(\phi(c(z))\bigr) = l^{\bot}\bigl(c(z)\bigr) = m\bigl(f(z)\bigr) = k\bigl(c'(z)\bigr)~,
\]
and $k$ is injective, so $c \fatsemi \phi = c'$.

\emph{Uniqueness.}
If $\psi$ also satisfies $\psi \fatsemi k = l^{\bot}$, then
$\psi \fatsemi k = \phi \fatsemi k$ and $k$ is monic, so $\psi = \phi$.
\end{proof}

\section{Congruence}
\label{sec:appendix-rewriting-congruence}

\begin{lemma}[Congruence]
Let $\mathfrak{S}$ be any set of EDPOI rewrite rules and suppose
$\mathcal{F} \;\overline{\rightsquigarrow}_{\mathfrak{S}}\; \mathcal{G}$.
Then, whenever the composites are defined,
\begin{enumerate}
  \item $\mathcal{E} \fatsemi \mathcal{F} \;\overline{\rightsquigarrow}_{\mathfrak{S}}\; \mathcal{E} \fatsemi \mathcal{G}$
        and $\mathcal{F} \fatsemi \mathcal{E} \;\overline{\rightsquigarrow}_{\mathfrak{S}}\; \mathcal{G} \fatsemi \mathcal{E}$;
  \item $\mathcal{E} \otimes \mathcal{F} \;\overline{\rightsquigarrow}_{\mathfrak{S}}\; \mathcal{E} \otimes \mathcal{G}$
        and $\mathcal{F} \otimes \mathcal{E} \;\overline{\rightsquigarrow}_{\mathfrak{S}}\; \mathcal{G} \otimes \mathcal{E}$;
  \item for $1 \leq r \leq p$ 
  \[\hspace{-2em}
  \mathcal{F}^{1} \join \cdots \join \mathcal{F}^{r-1} \join \mathcal{F} \join \mathcal{F}^{r+1} \join \cdots \join \mathcal{F}^{p}
        \;\overline{\rightsquigarrow}_{\mathfrak{S}}\;
        \mathcal{F}^{1} \join \cdots \join \mathcal{F}^{r-1} \join \mathcal{G} \join \mathcal{F}^{r+1} \join \cdots \join \mathcal{F}^{p}~.
        \]
\end{enumerate}
Consequently $\sim_{\mathfrak{S}}$ is a congruence for $\fatsemi$, $\otimes$ and
$\join$.
\end{lemma}
\begin{proof}
Let $\mathcal{F} \;\overline{\rightsquigarrow}_{\mathfrak{S}}\; \mathcal{G}$. 
By~\cref{def:dpoi_rooted} there is a rule 
\[
\mathcal{L} \xleftarrow{f_{\mathrm{int}}} i' \coprod j' \xleftarrow{f_{\mathrm{ext}}} i \coprod j \xrightarrow{p_{\mathrm{ext}}} i'' \coprod j'' \xrightarrow{p_{\mathrm{int}}} \mathcal{R}~,
\]
such that the following diagram commutes, satisfying the conditions of~\cref{def:dpoi_rooted}.
% https://q.uiver.app/#q=WzAsMTIsWzAsMCwiXFxtYXRoY2Fse0x9Il0sWzEsMCwiaScgXFxjb3Byb2QgaiciXSxbMiwwLCJpIFxcY29wcm9kIGoiXSxbMCwyLCJcXG1hdGhjYWx7Rn0iXSxbMiwyLCJcXG1hdGhjYWx7TH1ee1xcYm90fSJdLFsyLDUsIm4gXFxjb3Byb2QgayJdLFsxLDMsIm4nIFxcY29wcm9kIGsnIl0sWzQsMiwiXFxtYXRoY2Fse0d9Il0sWzQsMCwiXFxtYXRoY2Fse1J9Il0sWzMsMywibicnIFxcY29wcm9kIGsnJyJdLFszLDAsImknJ1xcY29wcm9kIGonJyJdLFsyLDMsIihuJyBcXHNldG1pbnVzIFNfbikgXFxjb3Byb2QgKGsgXFxzZXRtaW51cyBTX2spIl0sWzAsMywibSIsMl0sWzEsMCwiZl97XFxtYXRocm17aW50fX0iLDJdLFsyLDEsImZfe1xcbWF0aHJte2V4dH19IiwyXSxbMiw0LCJjIl0sWzQsMywibF57XFxib3R9IiwyXSxbNSw2LCJnX3tcXG1hdGhybXtleHR9fSJdLFs2LDMsImdfe1xcbWF0aHJte2ludH19Il0sWzMsMiwiIiwxLHsic3R5bGUiOnsibmFtZSI6ImNvcm5lciJ9fV0sWzQsNywiXFx3aWRldGlsZGV7XFxpb3RhX3tcXGJvdH19Il0sWzgsNywiXFx3aWRldGlsZGV7XFxpb3RhX3tcXG1hdGhjYWx7Un19fSJdLFs1LDksImhfe1xcbWF0aHJte2V4dH19IiwyXSxbOSw3LCJoX3tcXG1hdGhybXtpbnR9fSIsMl0sWzcsMiwiIiwxLHsic3R5bGUiOnsibmFtZSI6ImNvcm5lciJ9fV0sWzIsMTAsInBfe1xcbWF0aHJte2V4dH19Il0sWzEwLDgsInBfe1xcbWF0aHJte2ludH19Il0sWzUsMTEsIlxcb3ZlcmxpbmV7Z31fe1xcbWF0aHJte2V4dH19IiwyXSxbMTEsNCwiZ157XFxib3R9IiwyXV0=
\begin{equation}
  \begin{tikzcd}
	{\mathcal{L}} & {i' \coprod j'} & {i \coprod j} & {i''\coprod j''} & {\mathcal{R}} \\
	\\
	{\mathcal{F}} && {\mathcal{L}^{\bot}} && {\mathcal{G}} \\
	& {n' \coprod k'} & {(n' \setminus S_n) \coprod (k' \setminus S_k)} & {n'' \coprod k''} \\
	\\
	&& {n \coprod k}
	\arrow["m"', from=1-1, to=3-1]
	\arrow["{f_{\mathrm{int}}}"', from=1-2, to=1-1]
	\arrow["{f_{\mathrm{ext}}}"', from=1-3, to=1-2]
	\arrow["{p_{\mathrm{ext}}}", from=1-3, to=1-4]
	\arrow["c", from=1-3, to=3-3]
	\arrow["{p_{\mathrm{int}}}", from=1-4, to=1-5]
	\arrow["{\widetilde{\iota_{\mathcal{R}}}}", from=1-5, to=3-5]
	\arrow["\lrcorner"{anchor=center, pos=0.125, rotate=90}, draw=none, from=3-1, to=1-3]
	\arrow["{l^{\bot}}"', from=3-3, to=3-1]
	\arrow["{\widetilde{\iota_{\bot}}}", from=3-3, to=3-5]
	\arrow["\lrcorner"{anchor=center, pos=0.125, rotate=180}, draw=none, from=3-5, to=1-3]
	\arrow["{g_{\mathrm{int}}}", from=4-2, to=3-1]
	\arrow["{g^{\bot}}"', from=4-3, to=3-3]
	\arrow["{[h_{n}, h_{k}]}"', from=4-4, to=3-5]
	\arrow["{g_{\mathrm{ext}}}", from=6-3, to=4-2]
	\arrow["{\overline{g}_{\mathrm{ext}}}"', from=6-3, to=4-3]
	\arrow["{h^{n}_{\mathrm{ext}} \coprod h^{k}_{\mathrm{ext}}}"', from=6-3, to=4-4]
\end{tikzcd}~.
\tag{1}
\label{eq:extended_complement_diagram}
\end{equation}

\begin{enumerate}[parsep=\parskip,
    listparindent=\parindent]
\item We prove
\[
\mathcal E\fatsemi\mathcal F
\;\overline{\rightsquigarrow}_{\mathfrak S}\;
\mathcal E\fatsemi\mathcal G.
\]

Write the extended cospans representing \(\mathcal E\) and \(\mathcal F\) as
\[
n_1
\xrightarrow{f_{1,\mathrm{ext}}}
n_1'
\xrightarrow{f_{1,\mathrm{int}}}
\mathcal E
\xleftarrow{g_{1,\mathrm{int}}}
k_1'
\xleftarrow{g_{1,\mathrm{ext}}}
n
\]
and
\[
n
\xrightarrow{f_{2,\mathrm{ext}}}
n'
\xrightarrow{f_{2,\mathrm{int}}}
\mathcal F
\xleftarrow{g_{2,\mathrm{int}}}
k'
\xleftarrow{g_{2,\mathrm{ext}}}
k,
\]
so that $g_{\mathrm{ext}} = f_{2,\mathrm{ext}} \coprod g_{2,\mathrm{ext}}$ and $g_{\mathrm{int}} = [f_{2, \mathrm{int}}, g_{2, \mathrm{int}}]$.

Let
% https://q.uiver.app/#q=WzAsNCxbMCwwLCJrXzEiXSxbMiwwLCJcXG1hdGhjYWx7Rn0iXSxbMCwyLCJcXG1hdGhjYWx7RX0iXSxbMiwyLCJcXG1hdGhjYWx7RX0gXFxjb3Byb2Rfe2tfMX0gXFxtYXRoY2Fse0Z9Il0sWzAsMSwidiJdLFswLDIsInUiLDJdLFsyLDMsIlxcdmFyZXBzaWxvbiIsMl0sWzEsMywiXFxwaGkiXSxbMywwLCIiLDEseyJzdHlsZSI6eyJuYW1lIjoiY29ybmVyIn19XV0=
\begin{equation}
  \begin{tikzcd}
	{n} && {\mathcal{F}} \\
	\\
	{\mathcal{E}} && {\mathcal{E} \coprod_{n} \mathcal{F}}
	\arrow["f_{2,\mathrm{ext}} \fatsemi f_{2, \mathrm{int}}", from=1-1, to=1-3]
	\arrow["g_{1,\mathrm{ext}} \fatsemi g_{1,\mathrm{int}}"', from=1-1, to=3-1]
	\arrow["\phi", from=1-3, to=3-3]
	\arrow["\varepsilon"', from=3-1, to=3-3]
	\arrow["\lrcorner"{anchor=center, pos=0.125, rotate=180}, draw=none, from=3-3, to=1-1]
\end{tikzcd}
\tag{2}
\label{eq:pushout_square_e_circ_f}
\end{equation}
be the pushout computing the carrier of
\(\mathcal E\fatsemi\mathcal F\). 
The composite $\mathcal{E} \fatsemi \mathcal{F}$ is then given by the following extended cospan
\[
n_1 \xrightarrow{f_{1,\mathrm{ext}} \fatsemi \iota_{1}} n_1' \coprod (n' \setminus n) \xrightarrow{q_1} \mathcal{E} \coprod_{n} \mathcal{F} \xleftarrow{q_2} k' \coprod (k'_1 \setminus n) \xleftarrow{g_{2,\mathrm{ext}} \fatsemi \iota_{1}} k~,
\]
where 
\[
q_1 = [f_{1, \mathrm{int}} \fatsemi \varepsilon, (f_{2, \mathrm{int}} \fatsemi \phi) \restriction_{(n' \setminus n)}]
\]
and
\[
q_2 = [g_{2,\mathrm{int}} \fatsemi \phi, (g_{1, \mathrm{int}} \fatsemi \varepsilon) \restriction_{(k_1' \setminus n)}]~.
\]

Writing
\[
n \xrightarrow{\overline{g}^{n}_{\mathrm{ext}}} (n' \setminus S_n) \xrightarrow{g^{\bot}_{n}} \mathcal{L}^{\bot} \xleftarrow{g^{\bot}_{k}} (k' \setminus S_k) \xleftarrow{\overline{g}^{k}_{\mathrm{ext}}} k~,
\]
we similarly get a composite $\mathcal{E} \fatsemi \mathcal{L}^{\bot}$ given by the extended cospan
\begin{equation}
n_1 \xrightarrow{f_{1, \mathrm{ext}} \fatsemi \iota_{1}} n'_1 \coprod ((n' \setminus S_n) \setminus n) \xrightarrow{q_1^{\bot}} \mathcal{E} \coprod_{n} \mathcal{L}^{\bot} \xleftarrow{q_2^{\bot}} (k' \setminus S_k) \coprod (k'_1 \setminus n) \xleftarrow{\overline{g}^{k}_{\mathrm{ext}} \fatsemi \iota_{1}} k~,
\tag{3}
\label{eq:extended_complement_composite_cospan}
\end{equation}
with 
\[
q_1^{\bot} = [f_{1, \mathrm{int}} \fatsemi \varepsilon^{\bot}, (g_{n}^{\bot} \fatsemi \phi^{\bot}) \restriction_{((n' \setminus S_{n}) \setminus n)}]
\] 
and 
\[
q_2^{\bot} = [g_{k}^{\bot} \fatsemi \phi^{\bot}, (g_{1,\mathrm{int}} \fatsemi \varepsilon^{\bot}) \restriction_{((k_1' \setminus n))}]~,
\]
where $(\mathcal{E} \coprod_{n} \mathcal{L}^{\bot}, \phi^{\bot}, \varepsilon^{\bot})$ is the pushout 
% https://q.uiver.app/#q=WzAsNCxbMCwwLCJrXzEiXSxbMiwwLCJcXG1hdGhjYWx7TH1ee1xcYm90fSJdLFswLDIsIlxcbWF0aGNhbHtFfSJdLFsyLDIsIlxcbWF0aGNhbHtFfSBcXGNvcHJvZF97a18xfSBcXG1hdGhjYWx7Rn0iXSxbMCwxLCJnX3sxLFxcbWF0aHJte2V4dH19XFxmYXRzZW1pIGdfezF9XntcXGJvdH0iXSxbMCwyLCJnX3sxLFxcbWF0aHJte2V4dH19IFxcZmF0c2VtaSBnX3sxLFxcbWF0aHJte2ludH19IiwyXSxbMiwzLCJcXHZhcmVwc2lsb25ee1xcYm90fSIsMl0sWzEsMywiXFxwaGlee1xcYm90fSJdLFszLDAsIiIsMSx7InN0eWxlIjp7Im5hbWUiOiJjb3JuZXIifX1dXQ==
\begin{equation}
  \begin{tikzcd}
	{n} && {\mathcal{L}^{\bot}} \\
	\\
	{\mathcal{E}} && {\mathcal{E} \coprod_{n} \mathcal{L}^{\bot}}
	\arrow["{\overline{g}^{n}_{\mathrm{ext}}\fatsemi g_{n}^{\bot}}", from=1-1, to=1-3]
	\arrow["{g_{1,\mathrm{ext}} \fatsemi g_{1,\mathrm{int}}}"', from=1-1, to=3-1]
	\arrow["{\phi^{\bot}}", from=1-3, to=3-3]
	\arrow["{\varepsilon^{\bot}}"', from=3-1, to=3-3]
	\arrow["\lrcorner"{anchor=center, pos=0.125, rotate=180}, draw=none, from=3-3, to=1-1]
\end{tikzcd}
\tag{4}
\label{eq:pushout_square_e_circ_lbot}
\end{equation}

We claim that
\[
i\coprod j
\xrightarrow{c\fatsemi\phi^\bot}
\mathcal E\coprod_n\mathcal L^\bot
\xrightarrow{\id_{\mathcal{E}} \coprod l^{\bot}}
\mathcal E\coprod_n\mathcal F
\]
is an extended pushout complement of \(f\) and
\(m\fatsemi\phi\).

% https://q.uiver.app/#q=WzAsOCxbMCwxLCJuIl0sWzEsMSwiXFxtYXRoY2FsIExeXFxib3QiXSxbMiwxLCJcXG1hdGhjYWwgRiJdLFswLDIsIlxcbWF0aGNhbCBFIl0sWzEsMiwiXFxtYXRoY2FsIEVcXGNvcHJvZF9uXFxtYXRoY2FsIExeXFxib3QiXSxbMiwyLCJcXG1hdGhjYWwgRVxcY29wcm9kX25cXG1hdGhjYWwgRi4iXSxbMSwwLCJpIFxcY29wcm9kIGoiXSxbMiwwLCJcXG1hdGhjYWx7TH0iXSxbMCwxLCJcXG92ZXJsaW5lIGdfe1xcbWF0aHJte2V4dH19XntcXCxufVxuICAgICAgXFxmYXRzZW1pIGdfbl5cXGJvdCJdLFswLDMsImdfezEsXFxtYXRocm17ZXh0fX1cXGZhdHNlbWkgZ197MSxcXG1hdGhybXtpbnR9fSIsMl0sWzEsMiwibF5cXGJvdCJdLFsxLDQsIlxccGhpXlxcYm90Il0sWzIsNSwiXFxwaGkiXSxbMyw0LCJcXHZhcmVwc2lsb25eXFxib3QiLDJdLFs0LDUsIlxcaWRfe1xcbWF0aGNhbCBFfVxcY29wcm9kX24gbF5cXGJvdCIsMl0sWzQsMCwiIiwxLHsic3R5bGUiOnsibmFtZSI6ImNvcm5lciJ9fV0sWzUsMSwiIiwxLHsic3R5bGUiOnsibmFtZSI6ImNvcm5lciJ9fV0sWzEsMywiKDEpIiwxLHsic3R5bGUiOnsiYm9keSI6eyJuYW1lIjoibm9uZSJ9LCJoZWFkIjp7Im5hbWUiOiJub25lIn19fV0sWzIsNCwiKDIpIiwxLHsic3R5bGUiOnsiYm9keSI6eyJuYW1lIjoibm9uZSJ9LCJoZWFkIjp7Im5hbWUiOiJub25lIn19fV0sWzYsMSwiYyIsMl0sWzYsNywiZiJdLFs3LDIsIm0iXSxbMiw2LCIiLDIseyJzdHlsZSI6eyJuYW1lIjoiY29ybmVyIn19XSxbNywxLCIoMykiLDEseyJzdHlsZSI6eyJib2R5Ijp7Im5hbWUiOiJub25lIn0sImhlYWQiOnsibmFtZSI6Im5vbmUifX19XV0=
\begin{equation}
\begin{tikzcd}
	& {i \coprod j} & {\mathcal{L}} \\
	n & {\mathcal L^\bot} & {\mathcal F} \\
	{\mathcal E} & {\mathcal E\coprod_n\mathcal L^\bot} & {\mathcal E\coprod_n\mathcal F}
	\arrow["f", from=1-2, to=1-3]
	\arrow["c"', from=1-2, to=2-2]
	\arrow["{(3)}"{description}, draw=none, from=1-3, to=2-2]
	\arrow["m", from=1-3, to=2-3]
	\arrow["{\overline g_{\mathrm{ext}}^{\,n}
	      \fatsemi g_n^\bot}", from=2-1, to=2-2]
	\arrow["{g_{1,\mathrm{ext}}\fatsemi g_{1,\mathrm{int}}}"', from=2-1, to=3-1]
	\arrow["{l^\bot}", from=2-2, to=2-3]
	\arrow["{(1)}"{description}, draw=none, from=2-2, to=3-1]
	\arrow["{\phi^\bot}", from=2-2, to=3-2]
	\arrow["\lrcorner"{anchor=center, pos=0.125, rotate=180}, draw=none, from=2-3, to=1-2]
	\arrow["{(2)}"{description}, draw=none, from=2-3, to=3-2]
	\arrow["\phi", from=2-3, to=3-3]
	\arrow["{\varepsilon^\bot}"', from=3-1, to=3-2]
	\arrow["\lrcorner"{anchor=center, pos=0.125, rotate=180}, draw=none, from=3-2, to=2-1]
	\arrow["{\id_{\mathcal E}\coprod_n l^\bot}"', from=3-2, to=3-3]
	\arrow["\lrcorner"{anchor=center, pos=0.125, rotate=180}, draw=none, from=3-3, to=2-2]
\end{tikzcd}~.
\tag{5}
\label{eq:stacked_pushout_diagram}
\end{equation}
The square (1) is a pushout by the diagram~\eqref{eq:pushout_square_e_circ_lbot}.
From the extended pushout complement~\eqref{eq:extended_complement_diagram} we have that $\overline g_{\mathrm{ext}}^{n} \fatsemi g_{n}^{\bot} \fatsemi l^{\bot} = f_{2,\mathrm{ext}} \fatsemi f_{2, \mathrm{int}}$.
Consider maps $\varepsilon : \mathcal{E} \to \mathcal{E} \coprod_{n} \mathcal{F}$ and $l^{\bot} \fatsemi \phi : \mathcal{L}^{\bot} \to \mathcal{E} \coprod_{n} \mathcal{F}$, they agree on $n$ since
\begin{align*}
g_{1,\mathrm{ext}} \fatsemi g_{1,\mathrm{int}} \fatsemi \varepsilon &= f_{2,\mathrm{ext}} \fatsemi f_{2, \mathrm{int}} \fatsemi \phi\\
&= \overline{g}_{\mathrm{ext}}^{n} \fatsemi g_{n}^{\bot} \fatsemi l^{\bot} \fatsemi \phi~,
\end{align*}
and therefore the universal property of the pushout $\mathcal{E} \coprod_{n} \mathcal{L}^{\bot}$ gives us a unique map $\id_{\mathcal{E}} \coprod l^{\bot} : \mathcal{E} \coprod_{n} \mathcal{L}^{\bot} \to \mathcal{E} \coprod_{n} \mathcal{F}$ so that the square (2) commutes.
From the same universal property we have that $\varepsilon^{\bot} \fatsemi (\id_{\mathcal{E}} \coprod l^{\bot}) = \varepsilon$ and the outermost rectangle is a pushout by the diagram~\eqref{eq:pushout_square_e_circ_f}.
By the cancellation property of pushouts, the data above imply that the square (2) is a pushout.
Square (3) is a pushout given by the extended pushout complement~\eqref{eq:extended_complement_diagram} and by the pasting property of pushouts, the rectangle (3)--(2) is a pushout.
Both maps out of $n$ in the pushout defining
$\mathcal E\coprod_n\mathcal F$ are strict, so $\phi$ is strict.
Likewise,
\[
g_{1,\mathrm{ext}}\fatsemi g_{1,\mathrm{int}}
\qquad\text{and}\qquad
\overline g_{\mathrm{ext}}^{\,n}\fatsemi g_n^\bot
\]
are strict, so $\phi^\bot$ is strict.
Since $m$ and $c$ are anchored, it follows that
\[
m\fatsemi\phi
\qquad\text{and}\qquad
c\fatsemi\phi^\bot
\]
are anchored.
Together with the pushout property of the outer rectangle, this establishes
condition~(1) of~\cref{def:extended_complement}.

We next verify condition~(2). Since \(i\coprod j\) is pointed discrete and
\(f:i\coprod j\to\mathcal L\) is strict, the explicit construction of the
outer pushout identifies only the root and top-level interface vertices of
\(\mathcal L\) with elements of
\(\mathcal E\coprod_n\mathcal L^\bot\). In particular, it introduces no new
immediate children below the image of a non-root edge of \(\mathcal L\).

Hence, if \(e\in E_{\mathcal L}\) with
\(e\neq\top_{\mathcal L}\) and
\[
(m\fatsemi\phi)_E(e)
<^\mu_{\mathcal E\coprod_n\mathcal F}
x,
\]
then \(x=(m\fatsemi\phi)(y)\) for some
\(y\in V_{\mathcal L}\coprod E_{\mathcal L}\) satisfying
\[
e<^\mu_{\mathcal L}y.
\]
Consequently,
\[
x\in\operatorname{im}(m\fatsemi\phi),
\]
as required.

The same principle applies to condition (3); since $m \fatsemi \phi$ is a pushout coprojection over pointed discrete $i \coprod j$, the only vertices that get identified come from $i \coprod j$ and since $f$ is strict are necessarily top-level.

We next verify condition~(4).
The internal input and output interfaces of
\(\mathcal E\fatsemi\mathcal F\) are
\[
n_1'\coprod(n'\setminus n)
\qquad\text{and}\qquad
k'\coprod(k_1'\setminus n),
\]
respectively. According to condition~(4), the retained interfaces must
therefore be obtained by removing from these interfaces the affected
subinterfaces determined by the transported match \(m\fatsemi\phi\).
Under \(m\fatsemi\phi\), the images of the original subinterfaces
\(S_n\) and \(S_k\) lie in the \(n'\)- and \(k'\)-blocks of the above
interfaces, respectively.

Write \(S_{\mathrm{in}}\) and \(S_{\mathrm{out}}\) for these affected
subinterfaces, so that
\[
S_{\mathrm{in}}=\iota_2(S_n)
\qquad\text{and}\qquad
S_{\mathrm{out}}=\iota_1(S_k).
\]
Since \(\alpha_n\) and \(\alpha_k\) are canonical block isomorphisms,
\(S_n\) and \(S_k\) are unions of complete interface blocks. Hence, by
our convention on interface objects, omitting these blocks gives
\[
\begin{aligned}
\bigl(n_1'\coprod(n'\setminus n)\bigr)\setminus S_{\mathrm{in}}
&\cong
n_1'\coprod\bigl((n'\setminus S_n)\setminus n\bigr),\\
\bigl(k'\coprod(k_1'\setminus n)\bigr)\setminus S_{\mathrm{out}}
&\cong
(k'\setminus S_k)\coprod(k_1'\setminus n).
\end{aligned}
\]
These are precisely the internal interfaces of the previously
constructed composite extended cospan in
\eqref{eq:extended_complement_composite_cospan}. This retained extended
cospan is well-typed because it is the composite of the well-typed
extended cospan \(\mathcal E\) with the well-typed retained extended
cospan carried by \(\mathcal L^\bot\).

Let
\[
\jmath_{\mathrm{in}}:
n_1'\coprod\bigl((n'\setminus S_n)\setminus n\bigr)
\hookrightarrow
n_1'\coprod(n'\setminus n)
\]
and
\[
\jmath_{\mathrm{out}}:
(k'\setminus S_k)\coprod(k_1'\setminus n)
\hookrightarrow
k'\coprod(k_1'\setminus n)
\]
be the evident blockwise inclusions of the retained internal
interfaces. We need to check that
\[
\begin{aligned}
&
\Bigl(
(f_{1,\mathrm{ext}}\fatsemi\iota_1)
\coprod
(\overline g_{\mathrm{ext}}^{\,k}\fatsemi\iota_1)
\Bigr)
\fatsemi
(\jmath_{\mathrm{in}}\coprod\jmath_{\mathrm{out}})
\\
&\qquad =
(f_{1,\mathrm{ext}}\fatsemi\iota_1)
\coprod
(g_{2,\mathrm{ext}}\fatsemi\iota_1),
\end{aligned}
\]
and
\[
[q_1^\bot,q_2^\bot]
\fatsemi
(\id_{\mathcal E}\coprod l^\bot)
=
(\jmath_{\mathrm{in}}\coprod\jmath_{\mathrm{out}})
\fatsemi
[q_1,q_2].
\]
Both equalities may be verified on the respective coproduct blocks.

For the external legs, the input component is immediate because
\(\jmath_{\mathrm{in}}\) is the identity on the \(n_1'\)-block. On the
output component, we have
\[
(\overline g_{\mathrm{ext}}^{\,k}\fatsemi\iota_1)
\fatsemi\jmath_{\mathrm{out}}
=
g_{2,\mathrm{ext}}\fatsemi\iota_1,
\]
by condition~(4) for the original extended pushout complement, since
\(\jmath_{\mathrm{out}}\) restricts on the distinguished block to the
inclusion \(k'\setminus S_k\hookrightarrow k'\).

For the internal legs, consider the input component
\[
q_1^\bot
\fatsemi
(\id_{\mathcal E}\coprod l^\bot)
=
\jmath_{\mathrm{in}}\fatsemi q_1.
\]
Recall that
\[
q_1^\bot
=
\left[
f_{1,\mathrm{int}}\fatsemi\varepsilon^\bot,
\;
(g_n^\bot\fatsemi\phi^\bot)
\restriction_{(n'\setminus S_n)\setminus n}
\right]
\]
and
\[
q_1
=
\left[
f_{1,\mathrm{int}}\fatsemi\varepsilon,
\;
(f_{2,\mathrm{int}}\fatsemi\phi)
\restriction_{n'\setminus n}
\right].
\]
On the \(n_1'\)-block, the required equality is
\[
f_{1,\mathrm{int}}
\fatsemi\varepsilon^\bot
\fatsemi
(\id_{\mathcal E}\coprod l^\bot)
=
f_{1,\mathrm{int}}\fatsemi\varepsilon,
\]
which follows from the defining equations of
\(\id_{\mathcal E}\coprod l^\bot\).

On the retained distinguished block
\((n'\setminus S_n)\setminus n\), condition~(4) for the original
extended pushout complement gives
\[
g_n^\bot\fatsemi l^\bot
=
\jmath_n\fatsemi f_{2,\mathrm{int}},
\]
where
\[
\jmath_n:n'\setminus S_n\hookrightarrow n'.
\]
Together with
\[
\phi^\bot
\fatsemi
(\id_{\mathcal E}\coprod l^\bot)
=
l^\bot\fatsemi\phi,
\]
this gives
\[
\begin{aligned}
&
\left(
g_n^\bot\fatsemi\phi^\bot
\right)
\restriction_{(n'\setminus S_n)\setminus n}
\fatsemi
(\id_{\mathcal E}\coprod l^\bot)
\\
&\qquad =
\left(
\jmath_n\fatsemi f_{2,\mathrm{int}}\fatsemi\phi
\right)
\restriction_{(n'\setminus S_n)\setminus n}.
\end{aligned}
\]
By the interface-block convention, this is precisely the restriction
of
\[
\jmath_{\mathrm{in}}\fatsemi q_1
\]
to the same block. Thus
\[
q_1^\bot
\fatsemi
(\id_{\mathcal E}\coprod l^\bot)
=
\jmath_{\mathrm{in}}\fatsemi q_1.
\]
The corresponding equality for \(q_2^\bot\) and \(q_2\) follows by the
same argument on the output-interface blocks. Hence condition~(4)
holds.

It remains to check condition~(5).
The transported match \(m\fatsemi\phi\) has anchor
\[
a'\defeq\phi_E(a).
\]
Recall that the affected subinterfaces of the composite are
\[
S_{\mathrm{in}}=\iota_2(S_n)
\qquad\text{and}\qquad
S_{\mathrm{out}}=\iota_1(S_k).
\]
The maps required by condition~(5) are therefore
\[
\alpha_n\fatsemi\iota_2:
i'\setminus i\longrightarrow S_{\mathrm{in}}
\qquad\text{and}\qquad
\alpha_k\fatsemi\iota_1:
j'\setminus j\longrightarrow S_{\mathrm{out}},
\]
where \(\iota_2\) and \(\iota_1\) denote the corestrictions of the
coprojections to \(S_{\mathrm{in}}\) and \(S_{\mathrm{out}}\),
respectively.

Define the re-anchored interface maps
\[
\widehat q_1^{\,a'}
\defeq
\left\langle
a';
q_{1,V}\restriction_{V_{S_{\mathrm{in}}}}
\right\rangle
:
S_{\mathrm{in}}
\longrightarrow
\mathcal E\coprod_n\mathcal F
\]
and
\[
\widehat q_2^{\,a'}
\defeq
\left\langle
a';
q_{2,V}\restriction_{V_{S_{\mathrm{out}}}}
\right\rangle
:
S_{\mathrm{out}}
\longrightarrow
\mathcal E\coprod_n\mathcal F.
\]
It is sufficient to show that the following pasted diagrams commute:
\[
\begin{tikzcd}
	{i'\setminus i} && {\mathcal{L}} \\
	\\
	{S_n} && {\mathcal{F}} \\
	\\
	{S_{\mathrm{in}}} && {\mathcal{E}\coprod_{n}\mathcal{F}}
	\arrow["{f_{\mathrm{int}}\restriction}", from=1-1, to=1-3]
	\arrow["{\alpha_n}"', from=1-1, to=3-1]
	\arrow["m", from=1-3, to=3-3]
	\arrow["{\widehat{g}_{n}^{\,a}}"', from=3-1, to=3-3]
	\arrow["{\iota_{2}}"', from=3-1, to=5-1]
	\arrow["\phi"', from=3-3, to=5-3]
	\arrow["{\widehat{q}^{\,a'}_1}", from=5-1, to=5-3]
\end{tikzcd}
\qquad
\begin{tikzcd}
	{j'\setminus j} && {\mathcal{L}} \\
	\\
	{S_k} && {\mathcal{F}} \\
	\\
	{S_{\mathrm{out}}} && {\mathcal{E}\coprod_{n}\mathcal{F}}
	\arrow["{f_{\mathrm{int}}\restriction}", from=1-1, to=1-3]
	\arrow["{\alpha_k}"', from=1-1, to=3-1]
	\arrow["m", from=1-3, to=3-3]
	\arrow["{\widehat{g}_{k}^{\,a}}"', from=3-1, to=3-3]
	\arrow["{\iota_{1}}"', from=3-1, to=5-1]
	\arrow["\phi"', from=3-3, to=5-3]
	\arrow["{\widehat{q}^{\,a'}_2}", from=5-1, to=5-3]
\end{tikzcd}
\]
We will check the left diagram; the right diagram is analogous.
The upper square commutes by condition~(5) for the original extended
pushout complement. 
We check the lower square.

The maps $(\hat{g}_{n}^{\,a})_{E} \fatsemi \phi_{E}$ and $(\iota_2)_E \fatsemi (\widehat q_1^{\,a'})_E$ agree on the $\top_{S_n}$ by construction: $a' = \phi_{E}(a)$.
For every \(x\in V_{S_n}\), since \(S_n\subseteq n'\setminus n\) and
the second component of \(q_1\) is
\[
(f_{2,\mathrm{int}}\fatsemi\phi)
\restriction_{n'\setminus n},
\]
we have
\[
\begin{aligned}
\bigl(
\widehat g_n^{\,a}\fatsemi\phi
\bigr)_V(x)
&=
\phi_V\bigl(f_{2,\mathrm{int},V}(x)\bigr)
\\
&=
q_{1,V}\bigl(\iota_{2,V}(x)\bigr)
\\
&=
\bigl(
\iota_2\fatsemi\widehat q_1^{\,a'}
\bigr)_V(x).
\end{aligned}
\]
Thus
\[
\widehat g_n^{\,a}\fatsemi\phi
=
\iota_2\fatsemi\widehat q_1^{\,a'},
\]
and we get that the whole diagram commutes.

Since \(\alpha_n\) and \(\alpha_k\) are canonical block
isomorphisms, and the corestrictions
\[
\iota_2:S_n\longrightarrow S_{\mathrm{in}},
\qquad
\iota_1:S_k\longrightarrow S_{\mathrm{out}}
\]
are isomorphisms onto the corresponding blocks, the composites
\[
\alpha_n\fatsemi\iota_2
\qquad\text{and}\qquad
\alpha_k\fatsemi\iota_1
\]
are canonical isomorphisms of pointed discrete interfaces.
Hence we establish condition~(5).

We have now verified all five conditions of \cref{def:extended_complement}. 
Hence
\[
i\coprod j
\xrightarrow{c\fatsemi\phi^\bot}
\mathcal E\coprod_n\mathcal L^\bot
\xrightarrow{\id_{\mathcal E}\coprod l^\bot}
\mathcal E\coprod_n\mathcal F
\]
is an extended pushout complement of \(f\) and \(m\fatsemi\phi\).
It remains to show that completing this complement with the right-hand
pushout produces \(\mathcal E\fatsemi\mathcal G\).
First, recall the morphisms
\[
h^{n}_{\mathrm{ext}} : n \to n'', \qquad h_{n} : n'' \to \mathcal{G}
\]
\[
h_{\mathrm{ext}}^{k} : k \to k'', \qquad h_{k} : k'' \to \mathcal{G}
\]
We then have the following pushout diagram
% https://q.uiver.app/#q=WzAsNCxbMCwwLCJuIl0sWzIsMCwiXFxtYXRoY2Fse0d9Il0sWzAsMiwiXFxtYXRoY2Fse0V9Il0sWzIsMiwiXFxtYXRoY2Fse0V9IFxcY29wcm9kX3tufVxcbWF0aGNhbHtHfSJdLFswLDEsImhfe24sXFxtYXRocm17ZXh0fX1cXGZhdHNlbWkgaF97bicnLCBcXG1hdGhybXtpbnR9fSJdLFswLDIsImdfezEsXFxtYXRocm17ZXh0fX1cXGZhdHNlbWkgZ197MSxcXG1hdGhybXtpbnR9fSIsMl0sWzEsMywiXFxwc2kiXSxbMiwzLCJcXHZhcmVwc2lsb24nIiwyXSxbMywwLCIiLDEseyJzdHlsZSI6eyJuYW1lIjoiY29ybmVyIn19XV0=
\[\begin{tikzcd}
	n && {\mathcal{G}} \\
	\\
	{\mathcal{E}} && {\mathcal{E} \coprod_{n}\mathcal{G}}
	\arrow["{h^{n}_\mathrm{ext} \fatsemi h_{n}}", from=1-1, to=1-3]
	\arrow["{g_{1,\mathrm{ext}}\fatsemi g_{1,\mathrm{int}}}"', from=1-1, to=3-1]
	\arrow["\psi", from=1-3, to=3-3]
	\arrow["{\varepsilon'}"', from=3-1, to=3-3]
	\arrow["\lrcorner"{anchor=center, pos=0.125, rotate=180}, draw=none, from=3-3, to=1-1]
\end{tikzcd}\]
and the composite diagram below
% https://q.uiver.app/#q=WzAsOCxbMiwwLCJpIFxcY29wcm9kIGoiXSxbNCwwLCJcXG1hdGhjYWx7Un0iXSxbMiwyLCJcXG1hdGhjYWx7TH1ee1xcYm90fSJdLFs0LDIsIlxcbWF0aGNhbHtHfSJdLFsyLDQsIlxcbWF0aGNhbHtFfVxcY29wcm9kX3tufVxcbWF0aGNhbHtMfV57XFxib3R9Il0sWzQsNCwiXFxtYXRoY2Fse0V9XFxjb3Byb2RfblxcbWF0aGNhbHtHfSJdLFswLDIsIm4iXSxbMCw0LCJcXG1hdGhjYWx7RX0iXSxbMCwxLCJwIl0sWzEsMywiXFx3aWRldGlsZGV7XFxpb3RhX3tcXG1hdGhjYWx7Un19fSJdLFsyLDMsIlxcd2lkZXRpbGRle1xcaW90YV97XFxib3R9fSJdLFswLDIsImMiXSxbMiw0LCJcXHBoaV57XFxib3R9Il0sWzMsNSwiXFxwc2kiXSxbMywwLCIiLDIseyJzdHlsZSI6eyJuYW1lIjoiY29ybmVyIn19XSxbNSwyLCIiLDAseyJzdHlsZSI6eyJuYW1lIjoiY29ybmVyIn19XSxbNCw1LCJcXGlkX3tcXG1hdGhjYWx7RX19XFxjb3Byb2RcXHdpZGV0aWxkZXtcXGlvdGFfe1xcYm90fX0iLDJdLFs2LDIsIlxcb3ZlcmxpbmV7Z31eeyBufV97XFxtYXRocm17ZXh0fX0gXFxmYXRzZW1pIGdfe259XntcXGJvdH0iXSxbNiw3LCJnX3sxLFxcbWF0aHJte2V4dH19XFxmYXRzZW1pIGdfezEsXFxtYXRocm17aW50fX0iLDJdLFs3LDQsIlxcdmFyZXBzaWxvbl57XFxib3R9IiwyXSxbNCw2LCIiLDAseyJzdHlsZSI6eyJuYW1lIjoiY29ybmVyIn19XSxbMiw3LCIoMSkiLDEseyJzdHlsZSI6eyJib2R5Ijp7Im5hbWUiOiJub25lIn0sImhlYWQiOnsibmFtZSI6Im5vbmUifX19XSxbMyw0LCIoMikiLDEseyJzdHlsZSI6eyJib2R5Ijp7Im5hbWUiOiJub25lIn0sImhlYWQiOnsibmFtZSI6Im5vbmUifX19XSxbMSwyLCIoMykiLDEseyJzdHlsZSI6eyJib2R5Ijp7Im5hbWUiOiJub25lIn0sImhlYWQiOnsibmFtZSI6Im5vbmUifX19XV0=
\[\begin{tikzcd}
	&& {i \coprod j} && {\mathcal{R}} \\
	\\
	n && {\mathcal{L}^{\bot}} && {\mathcal{G}} \\
	\\
	{\mathcal{E}} && {\mathcal{E}\coprod_{n}\mathcal{L}^{\bot}} && {\mathcal{E}\coprod_n\mathcal{G}}
	\arrow["p", from=1-3, to=1-5]
	\arrow["c", from=1-3, to=3-3]
	\arrow["{(3)}"{description}, draw=none, from=1-5, to=3-3]
	\arrow["{\widetilde{\iota_{\mathcal{R}}}}", from=1-5, to=3-5]
	\arrow["{\overline{g}^{ n}_{\mathrm{ext}} \fatsemi g_{n}^{\bot}}", from=3-1, to=3-3]
	\arrow["{g_{1,\mathrm{ext}}\fatsemi g_{1,\mathrm{int}}}"', from=3-1, to=5-1]
	\arrow["{\widetilde{\iota_{\bot}}}", from=3-3, to=3-5]
	\arrow["{(1)}"{description}, draw=none, from=3-3, to=5-1]
	\arrow["{\phi^{\bot}}", from=3-3, to=5-3]
	\arrow["\lrcorner"{anchor=center, pos=0.125, rotate=180}, draw=none, from=3-5, to=1-3]
	\arrow["{(2)}"{description}, draw=none, from=3-5, to=5-3]
	\arrow["\psi", from=3-5, to=5-5]
	\arrow["{\varepsilon^{\bot}}"', from=5-1, to=5-3]
	\arrow["\lrcorner"{anchor=center, pos=0.125, rotate=180}, draw=none, from=5-3, to=3-1]
	\arrow["{\id_{\mathcal{E}}\coprod\widetilde{\iota_{\bot}}}"', from=5-3, to=5-5]
	\arrow["\lrcorner"{anchor=center, pos=0.125, rotate=180}, draw=none, from=5-5, to=3-3]
\end{tikzcd}\]
By a similar argument to the one before, 
\begin{align*}
\overline g_{\mathrm{ext}}^{\,n}
\fatsemi g_n^\bot
\fatsemi\widetilde\iota_{\bot}
\fatsemi\psi
&=
h_{\mathrm{ext}}^n
\fatsemi h_n
\fatsemi\psi
\\
&=
g_{1,\mathrm{ext}}
\fatsemi g_{1,\mathrm{int}}
\fatsemi\varepsilon',
\end{align*}
and $\id_{\mathcal{E}} \coprod \widetilde{\iota_{\bot}}$ with $\varepsilon^{\bot} \fatsemi (\id_{\mathcal{E}} \coprod \widetilde{\iota_{\bot}}) = \varepsilon'$ is the unique morphism induced by the universal property.
So the diagram (2) commutes and is a pushout by the cancellation property and the stacked rectangle (3)--(2) is a pushout by the pasting property.

The composite \(\mathcal E\fatsemi\mathcal G\) is represented by
\[
n_1
\xrightarrow{f_{1,\mathrm{ext}}\fatsemi\iota_1}
n_1'\coprod(n''\setminus n)
\xrightarrow{q_3}
\mathcal E\coprod_n\mathcal G
\xleftarrow{q_4}
k''\coprod(k_1'\setminus n)
\xleftarrow{h_{\mathrm{ext}}^k\fatsemi\iota_1}
k,
\]
where
\[
q_3
=
\left[
f_{1,\mathrm{int}}\fatsemi\varepsilon',
\;
(h_n\fatsemi\psi)\restriction_{n''\setminus n}
\right]
\]
and
\[
q_4
=
\left[
h_k\fatsemi\psi,
\;
(g_{1,\mathrm{int}}\fatsemi\varepsilon')
\restriction_{k_1'\setminus n}
\right].
\]

On the other hand, by the EDPOI rewriting definition, the internal interfaces of the
resulting extended cospan are
\[
\overline n''
\defeq
\left(
n_1'\coprod\bigl((n'\setminus S_n)\setminus n\bigr)
\right)
\coprod i''_\circ
\]
and
\[
\overline k''
\defeq
\left(
(k'\setminus S_k)\coprod(k_1'\setminus n)
\right)
\coprod j''_\circ,
\]
where
\[
i''_\circ=i''\setminus i,
\qquad
j''_\circ=j''\setminus j.
\]

Let
\[
\mu_n:
n_1'\coprod\bigl((n'\setminus S_n)\setminus n\bigr)
\longrightarrow
\overline n'',
\qquad
\mu_i:i''_\circ\longrightarrow\overline n''
\]
be the corresponding coprojections.

The new strictly internal input interface map contributed by the right-hand
side of the rule is
\[
\overline\rho_i
\defeq
\left\langle
\top_{\mathcal E\coprod_n\mathcal G};
\left(
p_{\mathrm{int}}
\fatsemi\widetilde\iota_{\mathcal R}
\fatsemi\psi
\right)_V
\restriction_{V_{i''_\circ}}
\right\rangle
:
i''_\circ
\longrightarrow
\mathcal E\coprod_n\mathcal G
\]
Equivalently, since \(\psi\) is strict,
\[
\overline\rho_i=\rho_i\fatsemi\psi~.
\]

The internal input leg produced by the EDPOI construction is therefore
\[
\overline h_{n_1}
\defeq
\left[
q_1^\bot
\fatsemi
\bigl(
\id_{\mathcal E}\coprod\widetilde\iota_\bot
\bigr),
\;
\overline\rho_i
\right]
:
\overline n''
\longrightarrow
\mathcal E\coprod_n\mathcal G
\]
The corresponding external input leg is
\[
\overline h_{\mathrm{ext}}^{\,n_1}
\defeq
f_{1,\mathrm{ext}}
\fatsemi\iota_1
\fatsemi\mu_n
:
n_1\longrightarrow\overline n'',
\]
where \(\iota_1\) is the first coprojection into the corresponding
retained interface. Likewise, we have the morphisms corresponding to
the external and internal output interfaces
\[
\overline h_{\mathrm{ext}}^{\,k}
:
k\longrightarrow\overline k''
\]
and
\[
\overline h_k
:
\overline k''
\longrightarrow
\mathcal E\coprod_n\mathcal G.
\]

It remains to show that the extended cospan representing the composite,
\[
n_1
\xrightarrow{f_{1,\mathrm{ext}}\fatsemi\iota_1}
n_1'\coprod(n''\setminus n)
\xrightarrow{q_3}
\mathcal E\coprod_n\mathcal G
\xleftarrow{q_4}
k''\coprod(k_1'\setminus n)
\xleftarrow{h_{\mathrm{ext}}^k\fatsemi\iota_1}
k,
\]
is isomorphic to the extended cospan produced by the EDPOI construction,
\[
n_1
\xrightarrow{\overline h_{\mathrm{ext}}^{\,n_1}}
\overline n''
\xrightarrow{\overline h_{n_1}}
\mathcal E\coprod_n\mathcal G
\xleftarrow{\overline h_k}
\overline k''
\xleftarrow{\overline h_{\mathrm{ext}}^{\,k}}
k.
\]
We consider the input side; the output side is analogous.

Recall that
\[
n''
=
(n'\setminus S_n)\coprod i''_\circ
\]
and
\[
\overline n''
=
\left(
n_1'\coprod\bigl((n'\setminus S_n)\setminus n\bigr)
\right)
\coprod i''_\circ.
\]
Since \(n\) lies in the \(n'\setminus S_n\)-block, the convention on
interface objects gives the canonical block isomorphisms
\[
\begin{aligned}
\overline n''
&=
\left(
n_1'\coprod\bigl((n'\setminus S_n)\setminus n\bigr)
\right)
\coprod i''_\circ
\\
&\cong
n_1'\coprod
\left(
\bigl((n'\setminus S_n)\setminus n\bigr)
\coprod i''_\circ
\right)
\\
&\cong
n_1'\coprod(n''\setminus n).
\end{aligned}
\]
Let
\[
\beta_n:
\overline n''
\longrightarrow
n_1'\coprod(n''\setminus n)
\]
denote the resulting canonical isomorphism. It is the identity on the
\(n_1'\)-block and sends the
\((n'\setminus S_n)\setminus n\)- and \(i''_\circ\)-blocks into the
\(n''\setminus n\)-block through \(\nu_n\) and \(\nu_i\),
respectively.

We need to check that
\[
f_{1,\mathrm{ext}}\fatsemi\iota_1
=
\overline h_{\mathrm{ext}}^{\,n_1}\fatsemi\beta_n.
\]
By definition of \(\overline h_{\mathrm{ext}}^{\,n_1}\), we have
\[
\begin{aligned}
\overline h_{\mathrm{ext}}^{\,n_1}\fatsemi\beta_n
&=
f_{1,\mathrm{ext}}
\fatsemi\iota_1
\fatsemi\mu_n
\fatsemi\beta_n
\\
&=
f_{1,\mathrm{ext}}\fatsemi\iota_1,
\end{aligned}
\]
because \(\mu_n\fatsemi\beta_n\) restricts on the \(n_1'\)-block to
the first coprojection into
\(n_1'\coprod(n''\setminus n)\).

It remains to verify compatibility with the internal input legs, that
is,
\[
\overline h_{n_1}
=
\beta_n\fatsemi q_3.
\]
Since
\[
\overline n''
=
\left(
n_1'\coprod\bigl((n'\setminus S_n)\setminus n\bigr)
\right)
\coprod i''_\circ,
\]
it is sufficient to check this equality on its three blocks.

On the \(n_1'\)-block, we have
\[
f_{1,\mathrm{int}}
\fatsemi\varepsilon^\bot
\fatsemi
\bigl(
\id_{\mathcal E}\coprod\widetilde\iota_\bot
\bigr)
=
f_{1,\mathrm{int}}\fatsemi\varepsilon',
\]
which is the first component of \(q_3\).

On the \((n'\setminus S_n)\setminus n\)-block, the defining equations
of the two pushouts and of \(h_n\) give
\[
\phi^\bot
\fatsemi
\bigl(
\id_{\mathcal E}\coprod\widetilde\iota_\bot
\bigr)
=
\widetilde\iota_\bot\fatsemi\psi
\]
and
\[
\nu_n\fatsemi h_n
=
g_n^\bot\fatsemi\widetilde\iota_\bot.
\]
Therefore,
\[
\begin{aligned}
\overline h_{n_1}
\restriction_{(n'\setminus S_n)\setminus n}
&=
\left(
g_n^\bot
\fatsemi\widetilde\iota_\bot
\fatsemi\psi
\right)
\restriction_{(n'\setminus S_n)\setminus n}
\\
&=
\left(
\nu_n\fatsemi h_n\fatsemi\psi
\right)
\restriction_{(n'\setminus S_n)\setminus n}
\\
&=
\left(
\beta_n\fatsemi q_3
\right)
\restriction_{(n'\setminus S_n)\setminus n},
\end{aligned}
\]
where the last equality follows because \(\beta_n\) sends this block
into \(n''\setminus n\) through \(\nu_n\), thereby selecting the
second component of \(q_3\).

Finally, on the \(i''_\circ\)-block, we have
\[
\overline\rho_i
=
\rho_i\fatsemi\psi
=
\nu_i\fatsemi h_n\fatsemi\psi.
\]
Since \(\beta_n\) sends \(i''_\circ\) into \(n''\setminus n\) through
\(\nu_i\), it again selects the second component of \(q_3\), and hence
\[
\begin{aligned}
\left(
\beta_n\fatsemi q_3
\right)
\restriction_{i''_\circ}
&=
\nu_i\fatsemi h_n\fatsemi\psi
\\
&=
\rho_i\fatsemi\psi
\\
&=
\overline\rho_i
\\
&=
\overline h_{n_1}\restriction_{i''_\circ}.
\end{aligned}
\]
Thus
\[
\overline h_{n_1}
=
\beta_n\fatsemi q_3.
\]

For the output side, the convention on interface objects similarly
gives a canonical block isomorphism
\[
\beta_k:
\overline k''
\longrightarrow
k''\coprod(k_1'\setminus n)
\]
which sends the \(k'\setminus S_k\)- and \(j''_\circ\)-blocks into
\(k''\) through \(\nu_k\) and \(\nu_j\), respectively, and sends the
\(k_1'\setminus n\)-block into the second summand. The same
componentwise calculation gives
\[
\overline h_{\mathrm{ext}}^{\,k}\fatsemi\beta_k
=
h_{\mathrm{ext}}^k\fatsemi\iota_1
\]
and
\[
\overline h_k
=
\beta_k\fatsemi q_4.
\]

The maps \(\beta_n\) and \(\beta_k\) are precisely the canonical
reassociations and permutations of complete interface blocks allowed
by our convention. Thus \(\beta_n\), the identity on
\(\mathcal E\coprod_n\mathcal G\), and \(\beta_k\) define an
isomorphism of extended cospans. Hence the right-hand pushout produces
a representative of \(\mathcal E\fatsemi\mathcal G\), as required, and
we have
\[
\mathcal E\fatsemi\mathcal F
\;\overline{\rightsquigarrow}_{\mathfrak S}\;
\mathcal E\fatsemi\mathcal G.
\]
The case
\[
\mathcal F\fatsemi\mathcal E
\;\overline{\rightsquigarrow}_{\mathfrak S}\;
\mathcal G\fatsemi\mathcal E
\]
is analogous.
\item The cases of $\mathcal{E} \otimes \mathcal{F} \;\overline{\rightsquigarrow}_{\mathfrak{S}}\; \mathcal{E} \otimes \mathcal{G}$ and $\mathcal{F} \otimes \mathcal{E} \;\overline{\rightsquigarrow}_{\mathfrak{S}}\; \mathcal{G} \otimes \mathcal{E}$ are also analogous to the above, noting that the tensor is given by a pushout.
\item We will prove $\mathcal{F}^{1} \join \cdots \join \mathcal{F} \join \cdots \mathcal{F}^{p} \;\overline{\rightsquigarrow}_{\mathfrak{S}}\; \mathcal{F}^{1} \join \cdots \join \mathcal{G} \join \cdots \join \mathcal{F}^{p}$.

Put
\[
\mathcal J_{\mathcal F}
\defeq
\mathcal F^1\join\cdots\join
\mathcal F^{r-1}\join\mathcal F\join
\mathcal F^{r+1}\join\cdots\join\mathcal F^p
\]
and
\[
\mathcal J_{\mathcal G}
\defeq
\mathcal F^1\join\cdots\join
\mathcal F^{r-1}\join\mathcal G\join
\mathcal F^{r+1}\join\cdots\join\mathcal F^p.
\]
Using the retained extended cospan carried by \(\mathcal L^\bot\), put
\[
\mathcal J_{\bot}
\defeq
\mathcal F^1\join\cdots\join
\mathcal F^{r-1}\join\mathcal L^\bot\join
\mathcal F^{r+1}\join\cdots\join\mathcal F^p.
\]

Let
\[
\jmath_{\mathcal F}:\mathcal F\longrightarrow\mathcal J_{\mathcal F}
\qquad\text{and}\qquad
\jmath_{\bot}:\mathcal L^\bot\longrightarrow\mathcal J_{\bot}
\]
denote the inclusions of the \(r\)-th components. The transported match is
\[
m^{\join}
\defeq
m\fatsemi\jmath_{\mathcal F}:
\mathcal L\longrightarrow\mathcal J_{\mathcal F},
\]
with anchor
\[
a^{\join}
\defeq
\jmath_{\mathcal F,E}(a).
\]

Choose the fresh root, the fresh \(\boxsort\)-edge, and the remaining
components of \(\mathcal J_{\bot}\) and \(\mathcal J_{\mathcal F}\)
compatibly. There is then a homomorphism
\[
l_{\join}^{\bot}:
\mathcal J_{\bot}\longrightarrow\mathcal J_{\mathcal F}
\]
which acts as \(l^\bot\) on the distinguished component
\(\mathcal L^\bot\), and as the identity on the fresh join structure and
all the other components. In particular,
\[
\jmath_{\bot}\fatsemi l_{\join}^{\bot}
=
l^\bot\fatsemi\jmath_{\mathcal F}.
\]

Define
\[
c^{\join}
\defeq
c\fatsemi\jmath_{\bot}:
i\coprod j\longrightarrow\mathcal J_{\bot}.
\]
We claim that
\[
i\coprod j
\xrightarrow{c^{\join}}
\mathcal J_{\bot}
\xrightarrow{l_{\join}^{\bot}}
\mathcal J_{\mathcal F}
\]
is an extended pushout complement of \(f\) and \(m^{\join}\).
Consider the following diagram
% https://q.uiver.app/#q=WzAsNixbMCwwLCJcXG1hdGhjYWx7TH0iXSxbMiwwLCJpIFxcY29wcm9kIGoiXSxbMCwyLCJcXG1hdGhjYWx7Rn0iXSxbMiwyLCJcXG1hdGhjYWx7TH1ee1xcYm90fSJdLFswLDQsIlxcbWF0aGNhbHtKfV97XFxtYXRoY2Fse0Z9fSJdLFsyLDQsIlxcbWF0aGNhbHtKfV97XFxib3R9Il0sWzAsMiwibSIsMl0sWzEsMCwiZiIsMl0sWzEsMywiYyJdLFszLDIsImxeXFxib3QiXSxbMiw0LCJcXGptYXRoX3tcXG1hdGhjYWx7Rn19IiwyXSxbMyw1LCJcXGptYXRoX3tcXGJvdH0iXSxbNSw0LCJsXntcXGJvdH1fe1xcam9pbn0iXSxbMiwxLCIiLDIseyJzdHlsZSI6eyJuYW1lIjoiY29ybmVyIn19XSxbNCwzLCIiLDAseyJzdHlsZSI6eyJuYW1lIjoiY29ybmVyIn19XV0=
\[\begin{tikzcd}
	{\mathcal{L}} && {i \coprod j} \\
	\\
	{\mathcal{F}} && {\mathcal{L}^{\bot}} \\
	\\
	{\mathcal{J}_{\mathcal{F}}} && {\mathcal{J}_{\bot}}
	\arrow["m"', from=1-1, to=3-1]
	\arrow["f"', from=1-3, to=1-1]
	\arrow["c", from=1-3, to=3-3]
	\arrow["\lrcorner"{anchor=center, pos=0.125, rotate=90}, draw=none, from=3-1, to=1-3]
	\arrow["{\jmath_{\mathcal{F}}}"', from=3-1, to=5-1]
	\arrow["{l^\bot}", from=3-3, to=3-1]
	\arrow["{\jmath_{\bot}}", from=3-3, to=5-3]
	\arrow["\lrcorner"{anchor=center, pos=0.125, rotate=90}, draw=none, from=5-1, to=3-3]
	\arrow["{l^{\bot}_{\join}}", from=5-3, to=5-1]
\end{tikzcd}\]
We will show that the bottom square is a pushout.
The square commutes by construction, so it remains to check the universal property.
Let $Q$ be a rooted e-hypergraph and let $j_1$ and $j_2$ be homomorphisms such that $l^{\bot} \fatsemi j_1 = \jmath_{\bot} \fatsemi j_2$.
We need to check that there is a unique map $w : \mathcal{J}_{\mathcal{F}} \longrightarrow Q$ such that
\[
\jmath_{\mathcal{F}} \fatsemi w = j_1 \qquad \text{and} \qquad l^{\bot}_{\join} \fatsemi w = j_2~.
\]
Define \(w\) on the image of \(\jmath_{\mathcal F}\) by
\[
w(\jmath_{\mathcal F}(x))\defeq j_1(x),
\]
and on the image of $l^{\bot}_{\join}$ by
\[
w(l^{\bot}_{\join}(y))\defeq j_2(y).
\]
These two images jointly contain every vertex and edge of
\(\mathcal J_{\mathcal F}\).

The definition is well-defined because \(l^{\bot}_{\join}\) is the identity
on the fresh join structure and on every component other than the
distinguished one, while on the image of \(\jmath_\bot\) it acts as
\(l^\bot\). On the overlap of the two images, the two prescriptions
agree by
\[
l^\bot\fatsemi j_1=\jmath_\bot\fatsemi j_2.
\]
Moreover, if
\[
l^{\bot}_{\join}(y)=l^{\bot}_{\join}(y'),
\]
the only non-trivial case is when \(y=\jmath_\bot(z)\) and
\(y'=\jmath_\bot(z')\) with \(l^\bot(z)=l^\bot(z')\). In that case,
\[
j_2(y)
=
j_1(l^\bot(z))
=
j_1(l^\bot(z'))
=
j_2(y'),
\]
so the value of \(w\) is independent of the chosen preimage.

Since \(j_1\) and \(j_2\) are homomorphisms and agree on the overlap,
\(w\) is a homomorphism. By construction,
\[
\jmath_{\mathcal F}\fatsemi w=j_1
\qquad\text{and}\qquad
l^{\bot}_{\join}\fatsemi w=j_2.
\]
The two images jointly cover \(\mathcal J_{\mathcal F}\), so these
equations also make \(w\) unique.
The upper square is a pushout by the original EDPOI square, and by the pasting property, the whole rectangle is a pushout.
By construction, $m^{\join}$ is anchored and therefore we successfully checked condition (1) of~\cref{def:extended_complement}.

We next verify condition~(2). Let
\[
e\in E_{\mathcal L},
\qquad
e\neq\top_{\mathcal L},
\]
and let
\[
x\in
V_{\mathcal J_{\mathcal F}}
\coprod
E_{\mathcal J_{\mathcal F}}
\]
satisfy
\[
(m^{\join})_{E}(e)
<^\mu_{\mathcal J_{\mathcal F}}
x.
\]
We must show that \(x\in\operatorname{im}(m^{\join})\).
By construction, $x$ has a preimage in $\mathcal{F}$, say $x = \jmath_{\mathcal{F}}(x')$, and $\jmath_{\mathcal{F}}$ reflects the immediate predecessor relation.
Therefore, $m_{E}(e) <^{\mu}_{\mathcal{F}} x'$ and by condition (2) of the original extended pushout complement, $x' \in \operatorname{im}(m)$ and hence $x = \jmath_{\mathcal{F}}(x') \in \operatorname{im}(m^{\join})$.

By construction, \(\jmath_{\mathcal F}\) is injective on vertices and
edges, while \(m\) is injective on the contents of each
\(\boxsort\)-edge of \(\mathcal L\).
Thus \(m^{\join} = m \fatsemi \jmath_{\mathcal F}\) is injective on the
contents of each \(\boxsort\)-edge of \(\mathcal L\), establishing
condition~(3).

We now check condition~(4).
Let
\[
N_{\mathcal F}
=
S\coprod n^1\coprod\cdots\coprod n'\coprod\cdots\coprod n^p
\]
and
\[
K_{\mathcal F}
=
T\coprod k^1\coprod\cdots\coprod k'\coprod\cdots\coprod k^p
\]
be the internal interfaces of \(\mathcal J_{\mathcal F}\), according
to~\cref{def:join_of_extended_cospans}.

Write
\[
S_{\mathrm{in}}\defeq\iota_n(S_n)
\qquad\text{and}\qquad
S_{\mathrm{out}}\defeq\iota_k(S_k)
\]
for the affected subinterfaces of \(N_{\mathcal F}\) and
\(K_{\mathcal F}\), where \(\iota_n\) and \(\iota_k\) are the
coprojections of the distinguished interface blocks. The transported
match \(m^{\join}\) affects only the distinguished component
\(\mathcal F\). Moreover, since \(\alpha_n\) and \(\alpha_k\) are
canonical block isomorphisms, \(S_{\mathrm{in}}\) and
\(S_{\mathrm{out}}\) are unions of complete interface blocks.
Therefore, by our convention on interface objects,
\[
\begin{aligned}
N_{\bot}
\defeq
N_{\mathcal F}\setminus S_{\mathrm{in}}
&\cong
S\coprod n^1\coprod\cdots
\coprod(n'\setminus S_n)
\coprod\cdots\coprod n^p,\\
K_{\bot}
\defeq
K_{\mathcal F}\setminus S_{\mathrm{out}}
&\cong
T\coprod k^1\coprod\cdots
\coprod(k'\setminus S_k)
\coprod\cdots\coprod k^p.
\end{aligned}
\]
This makes
\[
n
\xrightarrow{\mathcal J^n_{\bot,\mathrm{ext}}}
N_{\bot}
\xrightarrow{\mathcal J^n_{\bot,\mathrm{int}}}
\mathcal J_{\bot}
\xleftarrow{\mathcal J^k_{\bot,\mathrm{int}}}
K_{\bot}
\xleftarrow{\mathcal J^k_{\bot,\mathrm{ext}}}
k
\]
a well-typed extended cospan, since every component is well-typed.

Let
\[
\xi_n:N_{\bot}\hookrightarrow N_{\mathcal F}
\qquad\text{and}\qquad
\xi_k:K_{\bot}\hookrightarrow K_{\mathcal F}
\]
be the evident blockwise inclusions, and let
\[
n
\xrightarrow{\mathcal J^n_{\mathcal F,\mathrm{ext}}}
N_{\mathcal F}
\xrightarrow{\mathcal J^n_{\mathcal F,\mathrm{int}}}
\mathcal J_{\mathcal F}
\xleftarrow{\mathcal J^k_{\mathcal F,\mathrm{int}}}
K_{\mathcal F}
\xleftarrow{\mathcal J^k_{\mathcal F,\mathrm{ext}}}
k
\]
be the corresponding cospan for \(\mathcal J_{\mathcal F}\).

For the input interfaces, condition~(4) requires
\[
\mathcal J^n_{\bot,\mathrm{ext}}\fatsemi\xi_n
=
\mathcal J^n_{\mathcal F,\mathrm{ext}}
\]
and
\[
\mathcal J^n_{\bot,\mathrm{int}}\fatsemi l^\bot_{\join}
=
\xi_n\fatsemi\mathcal J^n_{\mathcal F,\mathrm{int}}.
\]
The first equation is immediate: both sides are the coprojection of
the unchanged external-interface block \(S\) into
\(N_{\mathcal F}\).

The second equation may be checked on the interface blocks. On \(S\)
and on every block \(n^q\neq n'\), the maps
\(\mathcal J^n_{\bot,\mathrm{int}}\) and
\(\mathcal J^n_{\mathcal F,\mathrm{int}}\) agree, while \(\xi_n\)
and \(l^\bot_{\join}\) act as the identity on the corresponding
blocks and their images.

On the retained distinguished block \(n'\setminus S_n\), we have
\[
\begin{aligned}
&
\left(
\mathcal J^n_{\bot,\mathrm{int}}
\fatsemi l^\bot_{\join}
\right)
\restriction_{n'\setminus S_n}
\\
&\qquad =
g_n^\bot
\fatsemi\jmath_{\bot}
\fatsemi l^\bot_{\join}
\\
&\qquad =
g_n^\bot
\fatsemi l^\bot
\fatsemi\jmath_{\mathcal F}
\\
&\qquad =
\jmath_n
\fatsemi g_{\mathrm{int}}
\fatsemi\jmath_{\mathcal F}
\\
&\qquad =
\left(
\xi_n
\fatsemi\mathcal J^n_{\mathcal F,\mathrm{int}}
\right)
\restriction_{n'\setminus S_n}.
\end{aligned}
\]
Here the second equality follows from
\[
\jmath_{\bot}\fatsemi l^\bot_{\join}
=
l^\bot\fatsemi\jmath_{\mathcal F},
\]
and the third is condition~(4) for the original extended pushout
complement. The two composites also agree on the root by strictness.
Thus
\[
\mathcal J^n_{\bot,\mathrm{int}}\fatsemi l^\bot_{\join}
=
\xi_n\fatsemi\mathcal J^n_{\mathcal F,\mathrm{int}}.
\]

The output-interface equations follow by the same blockwise argument,
using \(g_k^\bot\), \(\jmath_k\), and \(\xi_k\). Hence condition~(4)
holds.

It remains to check condition~(5).
The transported match \(m\fatsemi\jmath_{\mathcal F}\) has anchor
\[
a'\defeq\jmath_{\mathcal F,E}(a).
\]
Recall that the affected subinterfaces of \(\mathcal J_{\mathcal F}\)
are
\[
S_{\mathrm{in}}=\iota_n(S_n)
\qquad\text{and}\qquad
S_{\mathrm{out}}=\iota_k(S_k),
\]
where \(\iota_n\) and \(\iota_k\) denote the corestrictions of the
corresponding coprojections to these affected blocks. The maps required
by condition~(5) are therefore
\[
\alpha_n\fatsemi\iota_n:
i'\setminus i\longrightarrow S_{\mathrm{in}}
\qquad\text{and}\qquad
\alpha_k\fatsemi\iota_k:
j'\setminus j\longrightarrow S_{\mathrm{out}}.
\]
Since \(\alpha_n\) and \(\alpha_k\) are canonical block
isomorphisms, and \(\iota_n\) and \(\iota_k\) are isomorphisms onto
the corresponding affected blocks, these composites are canonical
isomorphisms of pointed discrete interfaces.

Define the re-anchored interface maps
\[
\widehat{(\mathcal J^n_{\mathcal F,\mathrm{int}})}^{\,a'}
\defeq
\left\langle
a';
\mathcal J^n_{\mathcal F,\mathrm{int},V}
\restriction_{V_{S_{\mathrm{in}}}}
\right\rangle
:
S_{\mathrm{in}}
\longrightarrow
\mathcal J_{\mathcal F}
\]
and
\[
\widehat{(\mathcal J^k_{\mathcal F,\mathrm{int}})}^{\,a'}
\defeq
\left\langle
a';
\mathcal J^k_{\mathcal F,\mathrm{int},V}
\restriction_{V_{S_{\mathrm{out}}}}
\right\rangle
:
S_{\mathrm{out}}
\longrightarrow
\mathcal J_{\mathcal F}.
\]
It is sufficient to show that the following pasted diagrams commute:
\[
\begin{tikzcd}
	{i'\setminus i} && {\mathcal L} \\
	\\
	{S_n} && {\mathcal F} \\
	\\
	{S_{\mathrm{in}}} && {\mathcal J_{\mathcal F}}
	\arrow["{f_{\mathrm{int}}\restriction}", from=1-1, to=1-3]
	\arrow["{\alpha_n}"', from=1-1, to=3-1]
	\arrow["m", from=1-3, to=3-3]
	\arrow["{\widehat g_n^{\,a}}"', from=3-1, to=3-3]
	\arrow["{\iota_n}"', from=3-1, to=5-1]
	\arrow["{\jmath_{\mathcal F}}"', from=3-3, to=5-3]
	\arrow[
	  "{\widehat{(\mathcal J^n_{\mathcal F,\mathrm{int}})}^{\,a'}}",
	  from=5-1,
	  to=5-3
	]
\end{tikzcd}
\qquad
\begin{tikzcd}
	{j'\setminus j} && {\mathcal L} \\
	\\
	{S_k} && {\mathcal F} \\
	\\
	{S_{\mathrm{out}}} && {\mathcal J_{\mathcal F}}
	\arrow["{f_{\mathrm{int}}\restriction}", from=1-1, to=1-3]
	\arrow["{\alpha_k}"', from=1-1, to=3-1]
	\arrow["m", from=1-3, to=3-3]
	\arrow["{\widehat g_k^{\,a}}"', from=3-1, to=3-3]
	\arrow["{\iota_k}"', from=3-1, to=5-1]
	\arrow["{\jmath_{\mathcal F}}"', from=3-3, to=5-3]
	\arrow[
	  "{\widehat{(\mathcal J^k_{\mathcal F,\mathrm{int}})}^{\,a'}}",
	  from=5-1,
	  to=5-3
	]
\end{tikzcd}
\]
We check the left diagram; the right diagram is analogous. The upper
square commutes by condition~(5) for the original extended pushout
complement. It remains to check the lower square.

On the root, both composites have value \(a'\), since
\[
\jmath_{\mathcal F,E}
\bigl(
\widehat g_{n,E}^{\,a}(\top_{S_n})
\bigr)
=
\jmath_{\mathcal F,E}(a)
=
a'
\]
and
\[
\widehat{(\mathcal J^n_{\mathcal F,\mathrm{int}})}^{\,a'}_E
\bigl(
\iota_{n,E}(\top_{S_n})
\bigr)
=
a'.
\]
For every \(x\in V_{S_n}\),
\[
\begin{aligned}
\widehat{(\mathcal J^n_{\mathcal F,\mathrm{int}})}^{\,a'}_V
\bigl(\iota_{n,V}(x)\bigr)
&=
\mathcal J^n_{\mathcal F,\mathrm{int},V}
\bigl(\iota_{n,V}(x)\bigr)
\\
&=
\jmath_{\mathcal F,V}
\bigl(g_{\mathrm{int},V}(x)\bigr)
\\
&=
\jmath_{\mathcal F,V}
\bigl(\widehat g_{n,V}^{\,a}(x)\bigr).
\end{aligned}
\]
Thus the lower square, and hence the whole left rectangle, commutes.
The right rectangle follows by the same argument. This proves
condition~(5).

Thus, we have proved that
\[
i\coprod j
\xrightarrow{c^{\join}}
\mathcal J_{\bot}
\xrightarrow{l^\bot_{\join}}
\mathcal J_{\mathcal F}
\]
is an extended pushout complement.

It remains to construct the second pushout and identify the resulting
extended cospan with \(\mathcal J_{\mathcal G}\).

Let
\[
\jmath_{\mathcal G}:
\mathcal G\longrightarrow\mathcal J_{\mathcal G}
\]
denote the inclusion of the distinguished component. Choosing the fresh
join structure compatibly, define
\[
\widetilde{\iota_{\bot}}^{\join}:
\mathcal J_{\bot}\longrightarrow\mathcal J_{\mathcal G}
\]
to act as
\(\widetilde{\iota_{\bot}}:\mathcal L^\bot\to\mathcal G\)
on the distinguished component and as the identity on the fresh join
structure and every other component. In particular,
\[
\jmath_{\bot}\fatsemi
\widetilde{\iota_{\bot}}^{\join}
=
\widetilde{\iota_{\bot}}\fatsemi
\jmath_{\mathcal G}.
\]

Consider the commutative diagram
\[
\begin{tikzcd}
	{i\coprod j} & {\mathcal R} \\
	{\mathcal L^\bot} & {\mathcal G} \\
	{\mathcal J_{\bot}} & {\mathcal J_{\mathcal G}}
	\arrow["p", from=1-1, to=1-2]
	\arrow["c"', from=1-1, to=2-1]
	\arrow["{\widetilde{\iota_{\mathcal R}}}", from=1-2, to=2-2]
	\arrow["{\widetilde{\iota_{\bot}}}"', from=2-1, to=2-2]
	\arrow["{\jmath_{\bot}}"', from=2-1, to=3-1]
	\arrow[
	  "\lrcorner"{anchor=center, pos=0.125, rotate=180},
	  draw=none,
	  from=2-2,
	  to=1-1
	]
	\arrow["{\jmath_{\mathcal G}}", from=2-2, to=3-2]
	\arrow[
	  "{\widetilde{\iota_{\bot}}^{\join}}"',
	  from=3-1,
	  to=3-2
	]
	\arrow[
	  "\lrcorner"{anchor=center, pos=0.125, rotate=180},
	  draw=none,
	  from=3-2,
	  to=2-1
	]
\end{tikzcd}
\]
The upper square is the second pushout square of the original EDPOI
rewrite. The lower square is a pushout by the same universal-property
argument used for the first pushout square. Therefore, the whole
rectangle is a pushout.

By definition, the extended cospan for \(\mathcal J_{\mathcal G}\) has
internal interfaces
\[
N_{\mathcal G}
=
S\coprod n^1\coprod\cdots\coprod n''\coprod\cdots\coprod n^p
\]
and
\[
K_{\mathcal G}
=
T\coprod k^1\coprod\cdots\coprod k''\coprod\cdots\coprod k^p,
\]
with the evident interface maps. The EDPOI construction, on the other
hand, gives the internal interfaces
\[
N'_{\mathcal G}
=
N_{\bot}\coprod i''_\circ
\qquad\text{and}\qquad
K'_{\mathcal G}
=
K_{\bot}\coprod j''_\circ.
\]
Since
\[
n''
=
(n'\setminus S_n)\coprod i''_\circ
\qquad\text{and}\qquad
k''
=
(k'\setminus S_k)\coprod j''_\circ,
\]
our convention on interface objects gives canonical block
isomorphisms
\[
\beta_n:
N'_{\mathcal G}\longrightarrow N_{\mathcal G}
\qquad\text{and}\qquad
\beta_k:
K'_{\mathcal G}\longrightarrow K_{\mathcal G}.
\]
They are the identity on the fresh join blocks and on every unaffected
component block, and reassociate the newly introduced blocks
\(i''_\circ\) and \(j''_\circ\) into the distinguished component.
Consequently, they are precisely the interface isomorphisms permitted
by our convention.

The external legs commute with \(\beta_n\) and \(\beta_k\) because
these isomorphisms are the identity on the fresh external-interface
blocks \(S\) and \(T\). For the internal legs, compatibility holds on
the fresh join structure and the unaffected components by construction.
On the retained parts of the distinguished component, it follows from
\[
h_n\restriction_{n'\setminus S_n}
=
g_n^\bot\fatsemi\widetilde{\iota_\bot},
\qquad
h_k\restriction_{k'\setminus S_k}
=
g_k^\bot\fatsemi\widetilde{\iota_\bot},
\]
and on the newly introduced parts it follows from
\[
h_n\restriction_{i''_\circ}
=
\rho_i,
\qquad
h_k\restriction_{j''_\circ}
=
\rho_j.
\]
The maps also agree on the roots by strictness. Therefore
\(\beta_n\), the identity on \(\mathcal J_{\mathcal G}\), and
\(\beta_k\) define an isomorphism of extended cospans.

Thus, we conclude that
\[
\mathcal J_{\mathcal F}
\;\overline{\rightsquigarrow}_{\mathfrak S}\;
\mathcal J_{\mathcal G}.
\]
\end{enumerate}
\end{proof}
